\chapter{Vacuum Circuit Analysis}


\begin{objectivebox}
\begin{itemize}
    \item Establish the topo-kinematic circuit identity: the vacuum as a fully characterisable transmission-line network with impedance $Z_0$, propagation constant $\gamma$, and reflection coefficient $\Gamma$.
    \item Derive the AVE Smith Chart as the primary engineering tool for vacuum impedance matching.
    \item Compute macroscopic non-linear topological stress directly from the second-order radiation pressure on a saturated dielectric lattice.
\end{itemize}
\end{objectivebox}

\section{The Topo-Kinematic Circuit Identity}
\label{sec:topo_kinematic}

The Vacuum Circuit Analysis (VCA) framework rests upon a dimensional isomorphism between continuum spatial mechanics and electrical network theory --- exact within the regime fence stated immediately below. This section states the mapping row by row, verifies it through independent thermodynamic consistency checks, and tabulates the complete translation dictionary.

% 🔴 SCOPE FENCE 2026-08-02 (manuscript-side propagation of def-tk1xfm, RATIFIED SOLID Grant 2026-07-21; Rule 12 KEEP-BOTH). The prior opening wording is PRESERVED in git and read: "rests upon a single, exact dimensional isomorphism between continuum spatial mechanics and electrical network theory. This section derives the mapping rigorously, ...". Two ratified qualifiers were missing from this section and are supplied below: (a) the REGIME FENCE (co-equality holds BELOW the band edge) and (b) the standing STRENGTH CEILING ("identity-by-translation, NOT emerges-from / NOT a derivation"). KB authority: ave-kb/common/vocabulary-register.md:435-441 (def-tk1xfm, adjudicated-meaning + status fields). NO row, number, equation or dimensional check in this section changes -- only the regime and the strength of the claim are scoped.
\noindent\textbf{Scope of the isomorphism (2026-08-02, per the ratified \texttt{def-tk1xfm} node; the section body below is preserved per Rule~12).} The Topo-Kinematic Isomorphism (Axiom~2) is ratified as an \textbf{exact, lossless, gain-1, pole-less, invertible co-equality of the mechanical and electrical descriptions \emph{below the band edge}} ($\omega\tau \ll 1$, the long-wave regime). The co-equality is \textbf{regime-scoped}: at band-edge scales the distributed-in-bond (arccos transmission-line) and lumped mass--spring descriptions book the propagation delay differently and are \emph{no longer co-equal}. Two consequences for the reading of this section. \textbf{(1)} The unqualified word ``exact'' is fenced --- it holds in the long-wave regime, not at every scale. \textbf{(2)} The \textbf{strength ceiling stands}: this is \emph{identity-by-translation}, \textbf{not} emerges-from and \textbf{not} a derivation of a mechanism. Where the text below says ``derived'' or ``proves'' (the six-row substitution and the work--energy / impedance consistency checks), read \emph{identity-by-translation verified for consistency}, not a derived mechanism; ratified-SOLID status means the node is canonical, not that it derives anything. The ceiling was reaffirmed by a computed null --- the Cleave-01 registry-pump Chern number returns $C = 0$ on both readings and both enantiomorphs (\texttt{clm-clvchn}), so no derived-mechanism instance has emerged. The dynamics (poles $=$ masses, $Q = 1/\alpha$ radiative leak, the Axiom-4 saturation trap) live in the resonator transfer function, \emph{not} in the Axiom-2 dictionary; conflating the two inflates a units bridge into a derived mechanism. Canonical: \kbleaf{ave-kb/common/vocabulary-register.md} (\texttt{def-tk1xfm}).

\subsection{Defining the Topological Conversion Constant}

The fundamental bridge between electrical and mechanical domains is the \textbf{Topological Conversion Constant} $\xi_{topo}$, defined as the ratio of the elementary charge to the spatial lattice pitch:
\begin{resultbox}{Topological Conversion Constant}
\begin{equation}
    \xi_{topo} \equiv \frac{e}{\ell_{node}} = \frac{1.602 \times 10^{-19} \text{ C}}{3.862 \times 10^{-13} \text{ m}} \approx 4.149 \times 10^{-7} \text{ C/m}
    \label{eq:xi_topo_vca}
\end{equation}
\end{resultbox}

This constant carries dimensions of $[\text{C/m}]$ and encodes the physical fact that, in a discrete lattice with pitch $\ell_{node}$, one quantum of electrical charge corresponds to one quantum of spatial displacement. Neither $e$ nor $\ell_{node}$ is adjustable; both are fixed by the lattice axioms (Axioms 1--2).

\subsection{Deriving the Six-Row Translation Table}

Each row of the Topo-Kinematic identity is derived by substituting $\xi_{topo}$ into the standard SI definition of the electrical quantity and isolating the resulting mechanical observable.

\paragraph{Row 1: Charge $\leftrightarrow$ Displacement.}
Electrical charge is the time integral of current. Substituting the kinematic mapping for current ($I \equiv \xi_{topo}\, v$):
\begin{resultbox}{Charge--Displacement Identity}
\begin{equation}
    Q = \int I \, dt = \xi_{topo} \int v \, dt = \xi_{topo}\, x
    \label{eq:charge_displacement}
\end{equation}
\end{resultbox}
\noindent Dimensional check: $[\text{C/m}] \times [\text{m}] = [\text{C}]$. \checkmark

\paragraph{Row 2: Current $\leftrightarrow$ Velocity.}
Differentiating the charge identity with respect to time:
\begin{resultbox}{Current--Velocity Identity}
\begin{equation}
    I = \frac{dQ}{dt} = \xi_{topo}\, \frac{dx}{dt} = \xi_{topo}\, v
    \label{eq:current_velocity}
\end{equation}
\end{resultbox}
\noindent The hardware velocity limit $v_{max} = c$ maps to a maximum circuit current:
\begin{equation}
    I_{max} = \xi_{topo}\, c = 4.149 \times 10^{-7} \times 2.998 \times 10^{8} \approx 124.4 \text{ A}
\end{equation}
This is the absolute slew-rate limit of the discrete lattice. No topological current can exceed $124.4$ A because no spatial velocity can exceed $c$.

\paragraph{Row 3: Voltage $\leftrightarrow$ Force.}
Voltage is the energy per unit charge ($V = W/Q$). Substituting $Q = \xi_{topo}\, x$ and $W = \int F\, dx$:
\begin{resultbox}{Voltage--Force Identity}
\begin{equation}
    V = \frac{W}{Q} = \frac{\int F\, dx}{\xi_{topo}\, x} \quad \Longrightarrow \quad V = \xi_{topo}^{-1}\, F
    \label{eq:voltage_force}
\end{equation}
\end{resultbox}
\noindent The maximum lattice force is the electromagnetic string tension $T_{EM} = m_e c^2 / \ell_{node} \approx 0.212$ N. The corresponding voltage is:
\begin{equation}
    V_{snap} = \xi_{topo}^{-1}\, T_{EM} = \frac{0.212}{4.149 \times 10^{-7}} \approx 511{,}000 \text{ V} = 511 \text{ kV}
\end{equation}
This recovers $V_{snap} = m_e c^2 / e$ exactly, confirming the identity.

\paragraph{Row 4: Inductance $\leftrightarrow$ Mass.}
Inductance opposes changes in current ($V = L\, dI/dt$). Substituting the voltage and current identities:
\begin{resultbox}{Inductance--Mass Identity}
\begin{equation}
    V = L \frac{dI}{dt} \quad \Longrightarrow \quad \xi_{topo}^{-1} F = L\, \xi_{topo} \frac{dv}{dt} = L\, \xi_{topo}\, a
    \quad \Longrightarrow \quad L = \xi_{topo}^{-2}\, m
    \label{eq:inductance_mass}
\end{equation}
\end{resultbox}
\noindent For the electron: $L_e = \xi_{topo}^{-2}\, m_e = m_e / \xi_{topo}^{2} \approx 5.292 \times 10^{-18}$ H ($\approx 5.3$ aH). This is the equivalent lumped inductance of a single electron's rest mass.

\paragraph{Row 5: Capacitance $\leftrightarrow$ Compliance.}
Capacitance stores charge per unit voltage ($C = Q/V$). Substituting:
\begin{resultbox}{Capacitance--Compliance Identity}
\begin{equation}
    C = \frac{Q}{V} = \frac{\xi_{topo}\, x}{\xi_{topo}^{-1}\, F} = \xi_{topo}^{2} \frac{x}{F} = \xi_{topo}^{2}\, \kappa
    \label{eq:capacitance_compliance}
\end{equation}
\end{resultbox}
\noindent where $\kappa = x/F = 1/k$ is the mechanical compliance (inverse spring constant). Dielectric breakdown occurs when the lattice displacement exceeds its absolute yield limit ($x > \ell_{node}$), which is structurally isomorphic to capacitor voltage exceeding the breakdown threshold.

\paragraph{Row 6: Resistance $\leftrightarrow$ Viscosity.}
Resistance dissipates power as $P = I^2 R$. Substituting:
\begin{resultbox}{Resistance--Viscosity Identity}
\begin{equation}
    R = \frac{V}{I} = \frac{\xi_{topo}^{-1} F}{\xi_{topo}\, v} = \xi_{topo}^{-2} \frac{F}{v} = \xi_{topo}^{-2}\, \eta
    \label{eq:resistance_viscosity}
\end{equation}
\end{resultbox}
\noindent where $\eta = F/v$ is the mechanical drag coefficient (units: kg/s). When the dielectric saturates and the vacuum enters the zero-impedance slipstream ($\eta \to 0$), the equivalent circuit resistance drops to zero---the hallmark of a superconductor.

\subsection{Complete Translation Dictionary}

Table~the main text consolidates the six identities into a single reference dictionary.

% ═══════════════════════════════════════════════════════════════
% DOMAIN TRANSLATION TABLE: Topo-Kinematic Circuit Identity
% Source: Extracted from Book 5 §12
% Canonical location: manuscript/common/translation_circuit.tex
% ═══════════════════════════════════════════════════════════════
%
% Usage: % ═══════════════════════════════════════════════════════════════
% DOMAIN TRANSLATION TABLE: Topo-Kinematic Circuit Identity
% Source: Extracted from Book 5 §12
% Canonical location: manuscript/common/translation_circuit.tex
% ═══════════════════════════════════════════════════════════════
%
% Usage: % ═══════════════════════════════════════════════════════════════
% DOMAIN TRANSLATION TABLE: Topo-Kinematic Circuit Identity
% Source: Extracted from Book 5 §12
% Canonical location: manuscript/common/translation_circuit.tex
% ═══════════════════════════════════════════════════════════════
%
% Usage: \input{../../common/translation_circuit.tex}
%        (from any chapter two levels deep in manuscript/)

\begin{table}[H]
\centering
\caption{Topo-Kinematic Circuit Identity.  Every row is derived from $\xi_{topo} = e/\ell_{node} \approx 4.149 \times 10^{-7}$ C/m (Axioms~1--2).}
\label{tab:trans_circuit}
\small
\renewcommand{\arraystretch}{1.3}
\begin{tabular}{@{}l l c l@{}}
\toprule
\textbf{Electrical} & \textbf{Mechanical} & \textbf{Mapping} & \textbf{Units Check} \\
\midrule
Charge $Q$        & Displacement $x$ & $Q = \xi\, x$             & $[\text{C/m}][\text{m}] = [\text{C}]$ \\
Current $I$       & Velocity $v$     & $I = \xi\, v$             & $[\text{C/m}][\text{m/s}] = [\text{A}]$ \\
Voltage $V$       & Force $F$        & $V = \xi^{-1} F$          & $[\text{m/C}][\text{N}] = [\text{V}]$ \\
Inductance $L$    & Mass $m$         & $L = \xi^{-2} m$          & $[\text{m}^2/\text{C}^2][\text{kg}] = [\text{H}]$ \\
Capacitance $C$   & Compliance $\kappa$ & $C = \xi^{2} \kappa$   & $[\text{C}^2/\text{m}^2][\text{m/N}] = [\text{F}]$ \\
Resistance $R$    & Viscosity $\eta$ & $R = \xi^{-2} \eta$       & $[\text{m}^2/\text{C}^2][\text{kg/s}] = [\Omega]$ \\
\bottomrule
\end{tabular}
\end{table}

%        (from any chapter two levels deep in manuscript/)

\begin{table}[H]
\centering
\caption{Topo-Kinematic Circuit Identity.  Every row is derived from $\xi_{topo} = e/\ell_{node} \approx 4.149 \times 10^{-7}$ C/m (Axioms~1--2).}
\label{tab:trans_circuit}
\small
\renewcommand{\arraystretch}{1.3}
\begin{tabular}{@{}l l c l@{}}
\toprule
\textbf{Electrical} & \textbf{Mechanical} & \textbf{Mapping} & \textbf{Units Check} \\
\midrule
Charge $Q$        & Displacement $x$ & $Q = \xi\, x$             & $[\text{C/m}][\text{m}] = [\text{C}]$ \\
Current $I$       & Velocity $v$     & $I = \xi\, v$             & $[\text{C/m}][\text{m/s}] = [\text{A}]$ \\
Voltage $V$       & Force $F$        & $V = \xi^{-1} F$          & $[\text{m/C}][\text{N}] = [\text{V}]$ \\
Inductance $L$    & Mass $m$         & $L = \xi^{-2} m$          & $[\text{m}^2/\text{C}^2][\text{kg}] = [\text{H}]$ \\
Capacitance $C$   & Compliance $\kappa$ & $C = \xi^{2} \kappa$   & $[\text{C}^2/\text{m}^2][\text{m/N}] = [\text{F}]$ \\
Resistance $R$    & Viscosity $\eta$ & $R = \xi^{-2} \eta$       & $[\text{m}^2/\text{C}^2][\text{kg/s}] = [\Omega]$ \\
\bottomrule
\end{tabular}
\end{table}

%        (from any chapter two levels deep in manuscript/)

\begin{table}[H]
\centering
\caption{Topo-Kinematic Circuit Identity.  Every row is derived from $\xi_{topo} = e/\ell_{node} \approx 4.149 \times 10^{-7}$ C/m (Axioms~1--2).}
\label{tab:trans_circuit}
\small
\renewcommand{\arraystretch}{1.3}
\begin{tabular}{@{}l l c l@{}}
\toprule
\textbf{Electrical} & \textbf{Mechanical} & \textbf{Mapping} & \textbf{Units Check} \\
\midrule
Charge $Q$        & Displacement $x$ & $Q = \xi\, x$             & $[\text{C/m}][\text{m}] = [\text{C}]$ \\
Current $I$       & Velocity $v$     & $I = \xi\, v$             & $[\text{C/m}][\text{m/s}] = [\text{A}]$ \\
Voltage $V$       & Force $F$        & $V = \xi^{-1} F$          & $[\text{m/C}][\text{N}] = [\text{V}]$ \\
Inductance $L$    & Mass $m$         & $L = \xi^{-2} m$          & $[\text{m}^2/\text{C}^2][\text{kg}] = [\text{H}]$ \\
Capacitance $C$   & Compliance $\kappa$ & $C = \xi^{2} \kappa$   & $[\text{C}^2/\text{m}^2][\text{m/N}] = [\text{F}]$ \\
Resistance $R$    & Viscosity $\eta$ & $R = \xi^{-2} \eta$       & $[\text{m}^2/\text{C}^2][\text{kg/s}] = [\Omega]$ \\
\bottomrule
\end{tabular}
\end{table}


\subsection{Substrate-Native Vocabulary --- Three-Column Rosetta-Stone Extension}
\label{sec:rosetta_stone_3col}

\textbf{Substrate-native vocabulary convention (canonical 2026-05-14).}
The translation dictionary above maps EE projection quantities to ME
projections. The universal substrate vocabulary (formalized in the Common
Foreword $\S$``Three Boundary Observables and the Substrate-Observability
Rule'' and \texttt{docs/glossary.md}) extends this to a third column ---
\textbf{substrate-native} --- naming what the substrate \emph{itself} carries
before any observation-frame projection. This subsection captures the canonical
three-column form. The translation discipline: \textbf{default to substrate-native
when writing about substrate dynamics; name the EE / ME projection only when
the math is committed to a specific observation frame.}

\begin{center}\small
\begin{tabular}{|p{2.8cm}|p{2.5cm}|p{2.5cm}|p{4.5cm}|}
\hline
\textbf{Substrate-native} & \textbf{EE projection} & \textbf{ME projection} & \textbf{Use when} \\
\hline
Substrate (the substrate) & vacuum / EM field & spacetime / elastic medium &
  Default: substrate. EE/ME when probe is specific. \\
\hline
Node & circuit node & lattice point &
  Default: substrate-vocabulary. The K4 4-port tetrahedral active site. \\
\hline
Bond & transmission line & spring &
  Default: bond. \emph{Spring} acceptable when invoking the substrate-scale Cosserat micropolar pedagogy (Vol 1 Ch~\ref{ch:macroscopic_moduli}). \\
\hline
State & voltage / $V_{\text{inc}}$ & force / displacement &
  Default: substrate. Specify projection when math is probe-coupled. \\
\hline
Propagation & current flow & --- &
  Default: substrate. How state evolves between nodes. \\
\hline
Impedance $Z$ & already universal & already universal & Use freely. \\
\hline
Saturation kernel $S(A) = \sqrt{1-A^2}$ & Schwinger / breakdown & yield / rupture &
  Use freely; Axiom~4 canonical. \\
\hline
Boundary & charged surface & yield surface &
  Default: substrate-vocabulary. Always. $\Gamma = -1$ saturation surface. \\
\hline
% [2026-08-02 ENVELOPE WORD-SENSE reconciliation to the ruled KB block. Prior row (git trail):
%  "Envelope & wavefunction support & zone of influence & Default: substrate. What the substrate
%  \emph{actually sees}." -- glossed "envelope" as a REGION, the pre-ruling sense. Authority:
%  ave-kb/common/envelope-anatomy.md, "Word-level ruling -- CONFIRMED (Grant in-chat 2026-07-15)":
%  the token "envelope" now canonically denotes the WALL -- surface (i), the fully-yielded
%  S(A)->0, |Gamma|=1 mirror; the balance shell (ii) and the knee / dress edge (iii) keep their
%  own names. Vocabulary correction only: no numerical value, radius or gate status is touched
%  (the leaf's radii stay gate-registered, NOT SOLID).]
Envelope & $|\Gamma| = 1$ mirror & yield surface &
  Ruled 2026-07-15: \emph{envelope} $=$ \textbf{the wall} --- the fully-yielded
  $S(A) \to 0$ surface. A surface, not a region. The \textbf{balance shell} and
  the \textbf{knee / dress edge} keep their own names
  (\kbleaf{ave-kb/common/envelope-anatomy.md}). \\
\hline
\textbf{Linking number $\mathcal{Q}$} & charge $Q$ & (no clean analog) &
  Integer per boundary; universal substrate; project to $Q$ for EE. \\
\hline
\textbf{Winding number $\mathcal{J}$} & spin $J$ / mag.~moment & rotation &
  Half-integer per SU(2) double-cover; universal substrate. \\
\hline
\textbf{Strain integral $\mathcal{M}$} & inductance $L$ / inertia (kg) & inertia &
  $\int_\Omega (n - 1)\, dV$; substrate mass-equivalent boundary observable. \\
\hline
Strain $A$ & voltage gradient & already universal & Default: substrate. \\
\hline
Machian & (no clean analog) & (no clean analog) & Use freely; relational substrate property. \\
\hline
\end{tabular}
\end{center}

\paragraph{The three substrate boundary invariants $\mathcal{M}, \mathcal{Q}, \mathcal{J}$
(canonical 2026-05-14 evening).} These three integers/half-integer are the
only externally observable invariants at any local $\Gamma = -1$ boundary
(substrate-observability rule, Common Foreword $\S$``Three Boundary Observables
and the Substrate-Observability Rule''). This is the no-hair theorem applied at
every scale: from outside a BH, observers see $M, Q, J$; from outside an electron,
observers see $m_e, e, \hbar/2$. Same structure, scale-invariant.

\begin{center}\small
\begin{tabular}{|p{1.3cm}|p{4.5cm}|p{6cm}|}
\hline
\textbf{Symbol} & \textbf{Canonical name} & \textbf{Operational definition} \\
\hline
$\mathcal{M}$ & Integrated strain integral & $\int_\Omega (n(\mathbf{r}) - 1)\, dV$; mass quantum $\hbar/(\ell_{\text{node}} c)$; per Eq.~\ref{eq:varactor} Virial sum recovery, $\mathcal{M} \cdot \ell_{\text{node}} \cdot c^2 = m_e c^2$ \\
\hline
$\mathcal{Q}$ & Boundary linking number & $\mathrm{Link}(\partial\Omega, \mathbf{F}_{\text{substrate}}) \in \mathbb{Z}$; one integer per primary flux unit \\
\hline
$\mathcal{J}$ & Boundary winding number & $\mathrm{Wind}(\partial\Omega)$; half-integer per SU(2) double-cover (spin-$\tfrac{1}{2}$ from $720°$ identity-closure mechanism) \\
\hline
\end{tabular}
\end{center}

\textbf{Engine implementation:} \texttt{src/ave/core/boundary\_invariants.py}
(13/13 tests PASS) provides reference implementations: \texttt{compute\_M},
\texttt{compute\_Q}, \texttt{compute\_J} for scalar substrate fields
(e.g., from MasterEquationFDTD).

\textbf{Cross-references:}
\begin{itemize}\itemsep 0pt
\item \textbf{Canonical statement:} Common Foreword $\S$``Three Boundary
  Observables and the Substrate-Observability Rule'' --- formal definitions,
  same-mechanism-at-all-scales claim, how to analyze any boundary in AVE.
\item \textbf{Cosmic-scale application:}
  $\mathcal{J}_{\text{cosmic}} = \Omega_{\text{freeze}} \cdot I_{\text{cosmic}}$
  encodes the cosmological initial condition (Vol~3 Ch~4 \S\ref{sec:tki_strain_snap}).
\item \textbf{Substrate vocabulary table:} \texttt{docs/glossary.md} (canonical
  substrate-native / EE / ME projection table).
\item \textbf{Universal Saturation-Kernel Catalog:} Backmatter
  Ch~\ref{app:universal_saturation_kernel} --- the A-034 catalog refines the
  cosmic-mechanism horizon (``God's Hand'') into a cosmic-parameter horizon vs.
  observable mechanism; the saturation-kernel mechanism IS observable at four
  smaller scales (BCS, solar flares, BH ring-down, Schwarzschild exact). % [2026-06-15 LF-03]: "NOAA solar flares" -> "solar flares". The mechanism-observability claim survives (flares are directly observed avalanche-snap events), but the NOAA-GOES validation framing is dropped here: per KB leaf solar-flares-led-avalanche.md the NOAA GOES comparison is a synthesized illustrative timeline, validation pending a live fetch.
\end{itemize}

\subsection{Self-Consistency Verification}

\paragraph{Work--Energy Theorem.}
The physical work done to charge a capacitor is $W = \int V\, dQ$. Substituting the Topo-Kinematic identities:
\begin{equation}
    W = \int (\xi_{topo}^{-1}\, F)(\xi_{topo}\, dx) = \int F\, dx
\end{equation}
The $\xi_{topo}$ factors cancel identically, recovering the mechanical work integral with no residual scaling. This proves the mapping is dimensionally exact and energy-conserving.

\paragraph{Impedance Cross-Check.}
The characteristic impedance of free space maps as:
\begin{equation}
    Z_{mech} = \xi_{topo}^{2}\, Z_0 = (4.149 \times 10^{-7})^2 \times 376.73 \approx 6.485 \times 10^{-11} \text{ kg/s}
\end{equation}
This is independently verifiable as the acoustic impedance of a single lattice node: $Z_{acoustic} = \rho_{bulk}\, c\, \ell_{node}^2 \approx 7.91 \times 10^6 \times 3.0 \times 10^8 \times (3.86 \times 10^{-13})^2 \approx 3.5 \times 10^{-10}$ kg/s. The factor-of-$\sim 5$ difference arises from the porosity correction ($p_c = 8\pi\alpha \approx 0.183$), confirming that the scaling is structurally consistent within the geometric packing fraction of the lattice.\section{Constitutive Circuit Models for Vacuum Non-Linearities}
\label{sec:vca_nonlinear}

Standard circuit simulators rely on ideal, linear RLC components. However, the physical substrate exhibits precise non-linear behaviors under extreme mechanical stress, each governed by the Axiom~4 saturation kernel $S(V) = \sqrt{1-(V/V_{yield})^2}$. This section derives the exact non-linear equivalent circuit component for each physical regime.

\subsection{The Metric Varactor --- Longitudinal Bond Compliance (Dielectric Yield)}

% 🔴 SECTOR RE-KEYING CORRECTION 2026-08-02 (Rule 12 KEEP-BOTH; manuscript-side propagation of a
%   ratified KB ruling. The printed prose, Eq. eq:varactor, its Taylor expansion and
%   Table tab:varactor_values below are PRESERVED VERBATIM -- no equation, no number and no
%   table row is edited.) KB twin: nonlinear-vacuum-capacitance.md:18 -- "Grade-fork RESOLVED = T2
%   (Grant 2026-06-30; def-vyvsn1 adjudicated ...)": the A1 varactor's saturation-completion
%   voltage is keyed on V_snap, NOT V_yield. Corrected constitutive equation at that leaf :27;
%   corrected table at :38-:46 (V/V_snap from sqrt(alpha)~0.085 to 1.000, i.e. 43.65 kV ->
%   511.0 kV); sector note at :14 (the printed identification is "the genesis-24 double-count");
%   heading retitled at :12. Engine already re-keyed: src/ave/solvers/spice_models/
%   ave_vacuum_cell.lib:61,:84 ("C_eff = C0 / S(V), knee at V_SNAP (A1 compliance)") and the KB
%   ch18 index :33. SCOPE: this note covers THIS subsection only; it does not adjudicate the
%   keying of any other section of this chapter.
\noindent\textbf{Sector re-keying correction (2026-08-02, \texttt{def-vyvsn1} grade-fork RESOLVED $=$ T2, Grant 2026-06-30; the subsection body below is preserved per Rule~12 KEEP-BOTH).} The text below keys the diverging compliance $C_{eff} = C_0/S$ on $V_{yield} = \sqrt{\alpha}\,V_{snap} \approx 43.65$~kV. That keying is \textbf{corrected}. \textbf{(1) Which reactance.} The diverging $C_0/S$ is the \textbf{longitudinal-A1 bond compliance} (the stretch-reactance $1/k_a$), \emph{not} the transverse dielectric capacitance; the transverse permittivity $\varepsilon_{eff} = \varepsilon_0 S$ and the across-gap cell capacitance an LCR meter reads, $C_{diel} \propto S$, \emph{roll off} instead. The two are \textbf{orthogonal} reactances (A1 $\perp$ T2) that share the EE name ``capacitance''; identifying them is the \emph{genesis-24 double-count}. \textbf{(2) Which voltage.} $V_{yield} = \sqrt{\alpha}\,V_{snap} \approx 43.65$~kV is the \textbf{transverse Cosserat ($T_2$) self-trap wall} --- the electron's confining $\Gamma = -1$ --- \emph{not} the A1 compliance bound. The longitudinal-A1 compliance diverges at the higher $V_{snap} = m_e c^2/e \approx 511$~kV, a factor $1/\sqrt{\alpha} \approx 11.7$ above $V_{yield}$. Read Eq.~\ref{eq:varactor}, its Taylor expansion and every row of Table~\ref{tab:varactor_values} with $V_{yield} \to V_{snap}$: the tabulated \emph{ratios} $C_{eff}/C_0$ and $S$ are unchanged, but the voltage column rescales by $1/\sqrt{\alpha}$ --- the printed $43.65$~kV endpoint is the $V/V_{snap} = \sqrt{\alpha} \approx 0.085$ \emph{first} row of the corrected table, not its divergence. \textbf{(3) Consequence.} A standing electron's A1 mass core operates at strain $A = V_{yield}/V_{snap} = \sqrt{\alpha} \approx 0.085$ --- deeply \textbf{sub-saturated} ($S \approx 0.996$) --- which is \emph{why} it binds: the $S \to 0$ varactor runaway never fires on the mass channel, and the confining wall is supplied by the transverse $T_2$ self-trap instead. Since $V_{yield} = \sqrt{\alpha}\,V_{snap}$ exactly, $A = \sqrt{\alpha}$ is an $\alpha$-echo operating point (Class-C), not a free parameter. Canonical: \kbleaf{ave-kb/vol4/circuit-theory/ch1-vacuum-circuit-analysis/nonlinear-vacuum-capacitance.md} (\texttt{def-vyvsn1}); the engine's SPICE cell is already keyed on \texttt{V\_SNAP}. \emph{Scope:} this note covers this subsection only.

As defined by Axiom~4, the effective compliance (capacitance) of the spatial substrate is bounded by the dielectric saturation limit at $V_{yield} = \sqrt{\alpha}\, V_{snap} \approx 43.65$ kV. The constitutive equation follows directly from the saturation kernel applied to the electric sector:
\begin{resultbox}{Vacuum Varactor Constitutive Equation}
\begin{equation}
    C_{eff}(V) = \frac{C_0}{\sqrt{1 - \left(\dfrac{V}{V_{yield}}\right)^{\!2}}} = \frac{C_0}{S(V)}
    \label{eq:varactor}
\end{equation}
\end{resultbox}

To verify consistency with the weak-field limit, the Taylor expansion about $V = 0$ yields:
\begin{equation}
    C_{eff}(V) = C_0 \left[1 + \frac{1}{2}\left(\frac{V}{V_{yield}}\right)^{\!2} + \frac{3}{8}\left(\frac{V}{V_{yield}}\right)^{\!4} + \cdots\right]
    \label{eq:varactor_taylor}
\end{equation}
At low voltages ($V \ll V_{yield}$), the leading correction is quadratic---identical to the Euler-Heisenberg effective Lagrangian of QED. The classical linear vacuum ($C_{eff} = C_0$) is recovered to arbitrary precision. Table~\ref{tab:varactor_values} evaluates the capacitance divergence at representative drive levels.

\medskip\noindent\textbf{The parity theorem --- the sub-yield difference tone is forbidden by inversion symmetry.} Because the varactor kernel is \emph{even} in the drive (Eq.~\ref{eq:varactor_taylor} carries only even powers of $V$), the restoring force is \emph{odd}, $F(r) = r + \tfrac12 r^3 + \cdots$: the reversible sub-yield vacuum is a pure-$\chi^3$, inversion-symmetric varactor. Under $V \to -V$ the equation of motion is odd, so a two-tone spectral line at $m\,\omega_{\text{lo}} + n\,\omega_{\text{hi}}$ exists \emph{iff} $m+n$ is odd. The literal difference tone $\omega_{\text{hi}} - \omega_{\text{lo}}$ is therefore \textbf{forbidden} below yield --- by inversion symmetry, not merely suppressed (measured $\sim\!3.4\times10^{-11}$ below the four-wave-mixing product) --- and the difference channel becomes an \emph{inversion-symmetry meter}: a planted even ($\chi^2$) term lights the forbidden tone monotonically, power $\propto\beta^2$ (fitted exponent $2.000$, $R^2 = 0.99999999$). The allowed four-wave-mixing beat ($2\omega_{\text{lo}} - \omega_{\text{hi}}$) scales as $A^6$ (clean-regime exponent $6.02$, $R^2 = 0.999999$; the steeper global $6.16$, $R^2 = 0.9997$, is $\chi^5$ stiffening at large amplitude --- the two $R^2$ values pair with \emph{different} exponents and must not be fused). This parity structure is a FORM-level consequence of the Axiom-4 kernel (no value imported); the canonical statement, its generative in-driver verification, and the \texttt{interface\_scoped\_\_BULK\_OPEN} scope of the \emph{frequency} vertex (the bulk four-photon vertex was \emph{not} probed --- the beat collapses $\approx$17.1 orders at bulk separation, tracking the frozen skin-suppression model) are at \kbleaf{ave-kb/common/universal-saturation-kernel-catalog.md} (\texttt{clm-invmtr}, \texttt{clm-a6chi3}; two-tone four-photon form factor, PR~\#610).

\begin{table}[H]
\centering
\caption{Effective vacuum capacitance vs.\ drive level. The divergence accelerates sharply above 90\% of $V_{yield}$.}
\label{tab:varactor_values}
\small
\begin{tabularx}{\textwidth}{@{}@{@{}}r r r r@{}}
\toprule
$V/V_{yield}$ & $V$ (kV) & $C_{eff}/C_0$ & $S(V)$ \\
\midrule
0.10  &  4.37  & 1.005 & 0.995 \\
0.50  & 21.83  & 1.155 & 0.866 \\
0.90  & 39.29  & 2.294 & 0.436 \\
0.99  & 43.21  & 7.089 & 0.141 \\
0.999 & 43.61  & 22.37 & 0.045 \\
1.000 & 43.65  & $\infty$ & 0  \\
\bottomrule
\end{tabularx}
\end{table}

\subsection{The Relativistic Inductor (Lorentz Saturation)}

Under the Topo-Kinematic identity, mass maps to inductance ($L = \xi_{topo}^{-2}\, m$) and velocity maps to current ($I = \xi_{topo}\, v$). The relativistic mass increase $m_{eff} = \gamma\, m_0$ therefore maps directly to a current-dependent inductance:
\begin{resultbox}{Relativistic Inductor}
\begin{equation}
    L_{eff}(I) = \frac{L_0}{\sqrt{1 - \left(\dfrac{I}{I_{max}}\right)^{\!2}}}, \qquad I_{max} = \xi_{topo}\, c \approx 124.4 \text{ A}
    \label{eq:relativistic_inductor}
\end{equation}
\end{resultbox}

This is structurally identical to the varactor equation (Eq.~\ref{eq:varactor}), with the substitution $V \to I$ and $V_{yield} \to I_{max}$. The symmetry is not coincidental: both are projections of the single Axiom~4 kernel onto the electric and magnetic sectors, respectively.

\paragraph{The magnetic-sector saturation primitive --- and the static-bias scope.}
The relativistic inductor is the \emph{magnetic-sector saturation primitive}: the kernel argument is the circulating current $I$, not the static magnetic-field magnitude $|\mathbf B|$. This keyed-argument distinction (the $\varepsilon$-varactor keys on the voltage $V$; the $\mu$-inductor keys on the current $I$) fixes the node's operating point under static external bias. A static external \textbf{electric} field is a genuine operating-point bias for the $V$-keyed varactor --- it loads the $\varepsilon$-sector ($S_\varepsilon < 1$), shifts the local index, and changes $Z$ (the vacuum-impedance-mirror / E-route, Ch.~\ref{ch:falsifiable_predictions}). A static external \textbf{magnetic} field, by contrast, is \emph{not} an operating point for the $I$-keyed inductor: with $\partial\mathbf B/\partial t = 0$ there is no induced internal vacuum circulation (Lenz), so $I_{\mathrm{vac}} = 0$, $S_\mu = 1$, $\mu_{eff} = \mu_0$, and the magnetic sector contributes \emph{no} index shift ($\delta n_\mu = 0$ exactly, independent of field strength). The substrate is transparent to a static $\mathbf B$; the only static-field route to vacuum birefringence is the electric (E-route). The full node-up small/large-signal derivation (the three operating regimes R1/R2/R3 and the static-field grade asymmetry) is canonical at \kbleaf{vol4/circuit-theory/ch1-vacuum-circuit-analysis/node-up-small-large-signal.md}, with the PVLAS/BMV static-B verdict (null \emph{consistent} with AVE) at \kbleaf{vol4/falsification/ch11-experimental-bench-falsification/pvlas-static-b-verdict.md}. ($B_{snap}$ enters as an \emph{energy-density} scale, not the inductor's kernel argument.)

\paragraph{Derivation of $E = mc^2$.}
The energy stored in a current-carrying inductor is $E = \tfrac{1}{2} L_{eff}\, I^2$. Expanding at low velocity ($v \ll c$):
\begin{equation}
    E = \frac{1}{2} \frac{L_0}{\sqrt{1-(v/c)^2}}\, (\xi_{topo}\, v)^2 = \frac{1}{2} \gamma\, \xi_{topo}^{-2}\, m_0\, \xi_{topo}^2\, v^2 = \frac{1}{2} \gamma\, m_0\, v^2
\end{equation}
At $v = 0$, the rest energy stored in the inductor's self-field is:
\begin{equation}
    E_0 = \tfrac{1}{2} L_0 I_{max}^2 = \tfrac{1}{2} (\xi_{topo}^{-2} m_0)(\xi_{topo} c)^2 = \tfrac{1}{2} m_0 c^2
\end{equation}
Summing with the capacitive (potential) energy via the Virial Theorem recovers $E_{total} = m_0 c^2$.

\paragraph{Why SPICE Cannot Exceed $c$.}
In a standard SPICE transient simulation, the inductor $L_0$ limits the current slew rate to $dI/dt = V/L$. As $I \to I_{max}$, $L_{eff} \to \infty$, and the slew rate collapses to zero. The circuit physically cannot push current (matter) past the hardware limit. SPICE engines natively enforce Special Relativity without any code modification.

\subsection{The Viscoelastic TVS Zener Diode (Phase Transition)}

% 🔴 SUB-YIELD BULK-LOSS FLAG 2026-08-02 (Rule 12 KEEP-BOTH; the TVS body, Eq. eq:tvs_transition,
%   the "highly resistive solid" sentence and the skin-effect subsection below are PRESERVED
%   VERBATIM -- no equation and no number is edited). Ruled authority:
%   ave-kb/common/substrate-native-terminology.md:33 (GREEN RULED 2026-07-19, Grant in-chat) --
%   "the corpus's genuine loss/irreversibility channels reduce from four to THREE: (1) radiative
%   port, (2) boundary-Joule extraction, (3) Regime-IV rupture"; a sub-yield BULK Re(Z) sink is
%   not among them (Ax 3). Term-register row :49 gives the substrate-native replacement for
%   "friction / drag": reflection / mismatch (Gamma != 0), with radiation resistance R_rad = Z_0
%   the real port. Consumer ruling: clm-h55fy1 at ave-kb/vol4/falsification/
%   ch11-experimental-bench-falsification/boundary-trapping-test.md:11 -- "the deep-space bulk
%   coupling to slow matter is lossless / pure-reactance", energy-shedding drag retracted.
%   Same-volume precedent banner: 11_experimental_falsification.tex:358 (2026-07-19). This is a
%   status flag, NOT a re-derivation: no replacement skin depth or drag coefficient is supplied.
\noindent\textbf{Scope flag (2026-08-02, sub-yield bulk loss; the body below is preserved per Rule~12 KEEP-BOTH).} This subsection and the skin-effect subsection that follows read the sub-yield vacuum as a \emph{dissipative} medium: ``a highly resistive solid \ldots kinematically gripping embedded matter via inductive drag ($\eta_0 > 0$)'', with $R_{vac} = \xi_{topo}^{-2}\eta_0 > 0$ giving a finite skin depth. That reading is \textbf{flagged}. Per the ruled loss-channel enumeration, the corpus's genuine loss/irreversibility channels are \textbf{three} --- the radiative port, boundary-Joule extraction at a resistive load, and Regime-IV rupture --- and a \emph{sub-yield bulk} $\mathrm{Re}(Z)$ sink is \textbf{not} among them (Axiom~3 forbids internal viscous heating in the reactive regime). The consumer ruling is explicit for slow matter in the bulk: the deep-space bulk coupling is \textbf{lossless / pure-reactance}, and the energy-shedding drag is retracted. What survives is the \textbf{reactive} reading: ``grip'' is impedance mismatch and reflection ($\Gamma \neq 0$), and the one honest resistive term is the \emph{radiation} resistance $R_{rad} \equiv Z_0$ (a real port, i.e.\ wave-making drag above the radiation threshold), not internal friction. Two riders. \textbf{(a)} The finite-skin-depth statement below \emph{consumes} the demoted $\eta_0$; no corrected skin depth is supplied here (that recomputation is not performed in this pass). \textbf{(b)} The Orbital Friction Paradox section later in this chapter resolves the paradox by \emph{neutralising} a presupposed sub-yield drag; the flag lands on that premise, not on that section's orthogonality conclusion ($P_{real} = F_g v_{orb}\cos 90^\circ = 0$), which is unaffected. Canonical: \kbleaf{ave-kb/common/substrate-native-terminology.md} (ruled 2026-07-19) and \kbleaf{ave-kb/vol4/falsification/ch11-experimental-bench-falsification/boundary-trapping-test.md} (\texttt{clm-h55fy1}); sibling banner in the \emph{Experimental Bench Falsification} chapter, \S``Orbital Detritus and Boundary Trapping'' (\texttt{sec:boundary\_trapping}).

Mutual inductance yields when the applied shear stress exceeds the structural threshold ($\tau > \tau_{yield}$). This voltage-driven breakdown maps exactly to a \textbf{Transient Voltage Suppression (TVS) Zener Diode}:

\begin{resultbox}{TVS Breakdown: Solid $\to$ Slipstream Transition}
\begin{equation}
    \eta_{eff}(V) = \begin{cases}
        \eta_0 & V < V_{yield} \quad \text{(solid: high drag)} \\[4pt]
        0      & V \geq V_{yield} \quad \text{(slipstream: zero drag)}
    \end{cases}
    \label{eq:tvs_transition}
\end{equation}
\end{resultbox}

Below the yield voltage, the vacuum acts as a highly resistive solid---kinematically gripping embedded matter via inductive drag ($\eta_0 > 0$). Above $V_{yield}$, an avalanche dielectric transition annihilates the mutual inductance. The vacuum enters ideal frictionless flow ($\eta = 0$): the \textbf{Zero-Impedance Slipstream}. The yield stress evaluates from the bulk energy density and packing geometry:
% [2026-07-19 Tier-2.5 tau_yield DISAMBIGUATION (KEEP-BOTH; NO renumber / NO redefinition / NO value change).]
%   WHICH tau_yield: this is the SHEAR / MUTUAL-INDUCTANCE branch, form
%   tau_yield = rho_bulk c^2 * (6 V_crossing) * p_c/8pi = rho_bulk c^2 * V_total * alpha ~ 1.04e22 Pa
%   (the shear stress at which the mutual inductance / inductive drag yields into the Zero-Impedance Slipstream).
%   It is ONE of THREE numerically distinct magnitudes the corpus labels "yield stress" under Bingham/plastic
%   naming, each from a DIFFERENT derivation:
%     (A) dimensional/Planck estimate  (hbar c / ell_node^4) * alpha^-2  ~ 7.21e34 Pa
%         (manuscript/vol_0_engineering_compendium/chapters/02_analytical_summaries.tex:17-36 -- and see the
%          resurrected-dropped-claim FLAG there vs backmatter/01_appendices.tex:77-88)
%     (B) THIS site  -- shear / mutual-inductance  rho_bulk c^2 V_total alpha  ~ 1.04e22 Pa  [here]
%     (C) magnetic-sector branch  tau_y = B_snap^2 / 2 mu_0
%         (manuscript/vol_1_foundations/chapters/06_universal_operators.tex:113)
%   Per-derivation / per-sector, not a single reconciled value; reconciliation ROUTED to the vol_0 compendium
%   reconciliation lane per research/2026-07-17_regime-iv-dissipation-audit.md:128 -- NOT adjudicated here.
\begin{align}
    \tau_{yield} &= \rho_{bulk}\, c^2 \times (6 \times \mathcal{V}_{crossing}) \times \frac{p_c}{8\pi} \nonumber \\
    &= \rho_{bulk}\, c^2 \times \mathcal{V}_{total} \times \alpha \nonumber \\
    &\approx 7.91 \times 10^{6} \times (3 \times 10^8)^2 \times 2.0 \times 0.00730 \nonumber \\
    &\approx 1.04 \times 10^{22}\ \text{Pa}
\end{align}
where $\mathcal{V}_{crossing} = V_{toroidal}/6 = 2.0/6$ is the per-crossing halo volume and $p_c/(8\pi) = \alpha \approx 1/137$ is the lattice porosity.

\subsection{The Vacuum Memristor (Thixotropic Hysteresis)}

% 🔴 FORK-OPEN SCOPE CAVEAT 2026-07-19 (yield-fork discriminators, Rule 12 KEEP-BOTH): the near-yield "dissipative / plastic / memristive hysteresis-loss" body of this subsection is PRESERVED verbatim; the caveat below scopes it as ONE branch of the registered-OPEN near-yield yield fork (finite-area memristive loop = dissipative vs zero-area saturating reactance = lossless refusal), NOT a settled result. Ground truth: research/2026-07-19_yield-fork-discriminators_result.md (Leg A verdict B excludes the two-tau thixotropic-rectifier sub-branch by derivation; Leg B NEITHER; crux relocates to #59 Flag F, an upstream derivation question unreached by either driver -- the fork stays OPEN; Grant leans reversible/reactive); ave-kb/common/retention-transition-split.md (the near-yield loop is fork-gated, fork held OPEN: "do not bank it either way in load-bearing prose"; that leaf's RULED-2026-07-19 stamp is the Op3-transduction ruling, not a fork adjudication). FLAG-DON'T-FIX embedded (loss-location tension vs backmatter/06 + #735 Leg-B); see _orchestration/2026-07-19_manuscript-cleanup-ledger.md C2-FLAG-1.
\noindent\textbf{Scope caveat (2026-07-19, yield-fork discriminators; the dissipative near-yield body below is preserved per Rule~12 KEEP-BOTH).} The near-yield ``dissipative / plastic / memristive hysteresis-loss'' picture of this subsection --- the finite-area pinched loop, the memristance $M$ in $\Omega$, and the ``energy dissipated per thixotropic yield--heal cycle'' --- is \textbf{one branch of a registered-OPEN fork}, not a settled result. The near-yield crossing is either \emph{first-order overdamped} (dissipative: a finite-area memristive loop, the reading below) or \emph{second-order reactive} (lossless: a zero-area saturating reactance, $I_S \neq 0$); which one holds is the upstream derivation question \textbf{Flag~F} (registered in \#59~\S12; derivation home: \kbleaf{tau-relax-derivation.md}), and it is \textbf{OPEN} (Grant leans reversible/reactive). The registered discriminators (\texttt{research/2026-07-19\_yield-fork-discriminators\_result.md}) \emph{excluded the two-$\tau$ thixotropic-rectifier sub-branch by derivation} (Leg~A verdict~B: the canonical kernel carries no genuine $\operatorname{sign}(dr/dt)$ memory) but did \textbf{not} close the fork against the lossless branch. Per canon (\kbleaf{ave-kb/common/retention-transition-split.md}; fork held OPEN), \emph{do not bank the loop either way in load-bearing prose}. \emph{Flag (unresolved, routed to Grant):} whether the loss maximises at $f \ll 1/\tau_{relax}$ (as stated at the end of this subsection), is zero-area/elastic there (as \texttt{backmatter/06\_spice\_verification\_manual.tex} states), or peaks at $\omega\tau \approx 0.9$ ($f\tau \sim 1$, the \#735 Leg-B Debye shape) is itself part of the open fork --- surfaced, not resolved.

In the near-yield dissipative channel (at/above $V_{yield}$, distinct from the sub-yield lossless-reactive regime), the dielectric saturation--plastic transition requires a finite geometric relaxation time to physically liquefy the lattice:
\begin{resultbox}{Thixotropic Relaxation Time}
\begin{equation}
    \tau_{relax} = \frac{\ell_{node}}{c} = \frac{3.862 \times 10^{-13}}{2.998 \times 10^{8}} \approx 1.288 \times 10^{-21} \text{ s}
    \label{eq:relaxation_time}
\end{equation}
\end{resultbox}

Because the vacuum cannot alter its inductive resistance instantaneously, its present-state impedance depends on the historical integral of applied stress. This is the defining characteristic of a \textbf{Memristor}: a circuit element whose resistance is a function of the cumulative charge that has passed through it.

The constitutive relation is:
\begin{equation}
    M(q) = \frac{d\Phi}{dq}, \qquad \Phi(t) = \int_{-\infty}^{t} V(\tau)\, d\tau
    \label{eq:memristor_constitutive}
\end{equation}
where $M$ is the memristance (units: $\Omega$) and $\Phi$ is the magnetic flux linkage. Under high-frequency AC topological stress, the memristive vacuum produces a classic \textbf{Pinched Hysteresis Loop}: the $V$--$I$ Lissajous figure passes through the origin but encloses a finite area proportional to the energy dissipated during each thixotropic yield--heal cycle.

At drive frequencies $f \gg 1/\tau_{relax} \approx 7.8 \times 10^{20}$ Hz, the vacuum responds too slowly to yield and behaves as a purely elastic (linear) medium. At $f \ll 1/\tau_{relax}$, complete yield and recovery occur within each cycle, producing maximum hysteresis loss. The crossover frequency is set entirely by the lattice pitch and $c$.

% 🔴 LOSS-LOCATION CORRECTION 2026-07-20 (derivation-backed, Rule 12 KEEP-BOTH): the sentence above ("At f<<1/tau, ... producing maximum hysteresis loss") is PRESERVED verbatim; the correction below is now DERIVATION-BACKED, sharpening the 2026-07-19 fork-open caveat's loss-location flag. Ground truth: research/2026-07-20_jomega-derivation_result.md (merged #751), §7 loss-location adjudication.
\noindent\textbf{Loss-location correction (2026-07-20, J($\omega$)-derivation-backed; the ``maximum hysteresis loss at $f \ll 1/\tau_{relax}$'' sentence above is preserved per Rule~12 KEEP-BOTH).} That sentence is the \textbf{excluded world-(c) rate-independent-plastic reading}. Under the Axiom-3-lossless z=3 bond-coupled bath, the spectral density $J(f \ll 1/\tau_{relax}) \to 0$ (both coupling models), so the quasi-static limit is \textbf{elastic / zero-area}, i.e.\ loss is \textbf{minimal}, not maximal, at low $f$ --- the \texttt{backmatter/06\_spice\_verification\_manual.tex} ``zero enclosed area'' statement is \emph{correct} in exactly this limit. The dissipative/reactive loop-area actually \emph{peaks} near $\omega\tau_{relax} \sim 1$ (the Debye/response rate, \#735 Leg-B (V,I) peak $\approx 0.911$); and even there, treating Eq.~2.1's first-order Debye peak as a genuine \emph{bath} dissipation requires $\omega\tau_{relax} \ll 1$ AND $\gamma_0 > 0$ (Ohmic coupling) --- so the peak sits \textbf{outside Eq.~2.1's own regime of validity as a bath reduction} (\texttt{research/2026-07-20\_jomega-derivation\_result.md} \S3, \S7; under super-Ohmic coupling $\gamma_0 = 0$ and Eq.~2.1 is not recoverable at all). The fork on whether the residual finite loop is dissipative (first-order overdamped) or lossless (second-order reactive) stays \textbf{OPEN} at \#59 Flag~F (Grant leans reversible/reactive) --- surfaced, not resolved.

\subsection{The Zero-Impedance Skin Effect (Metric Faraday Cages)}

In standard electrical engineering, the AC skin depth $\delta$ governs how deeply alternating current penetrates into a conductor:
\begin{resultbox}{Classical Skin Depth}
\begin{equation}
    \delta = \sqrt{\frac{2\rho}{\omega\mu}} = \sqrt{\frac{2 R_{vac}}{\omega\mu_0}}
    \label{eq:skin_depth}
\end{equation}
\end{resultbox}

Under the Topo-Kinematic identity, $R_{vac} \equiv \xi_{topo}^{-2}\, \eta_{vac}$: the vacuum's electrical resistance maps to its mutual inductance (drag coefficient). Evaluating the skin depth in the two vacuum phases:

\begin{enumerate}
    \item \textbf{Unsaturated Vacuum ($V < V_{yield}$):} $R_{vac} = \xi_{topo}^{-2}\, \eta_0 > 0$. The skin depth is finite; shear waves penetrate to a depth proportional to $\sqrt{\eta_0}$. At 1~GHz:
    \begin{equation}
        \delta_{solid} = \sqrt{\frac{2\, \xi_{topo}^{-2}\, \eta_0}{2\pi \times 10^9 \times \mu_0}} \sim \text{finite (deep-space penetration)}
    \end{equation}

    \item \textbf{Saturated Slipstream ($V \geq V_{yield}$):} $\eta_{eff} \to 0$, therefore $R_{vac} \to 0$:
    \begin{equation}
        \delta_{slipstream} = \sqrt{\frac{2 \times 0}{\omega\mu_0}} \equiv 0
    \end{equation}
\end{enumerate}

When the metric exceeds the yield threshold, the skin depth collapses to zero. The destructive, high-shear slipstream is confined \textit{entirely} to the exterior surface of the macroscopic body. The interior metric acts as a \textbf{Topological Faraday Cage}: perfectly shielded from external structural shear, even at extreme gravitational gradients. This provides the mechanical basis for why planetary interiors remain structurally intact inside deep gravity wells.
\section{The Impedance of Free Space (\texorpdfstring{$Z_0$}{Z0})}
\label{sec:z0_derivation}

A foundational parameter in classical electromagnetism is the Characteristic Impedance of Free Space, $Z_0 = \sqrt{\mu_0/\epsilon_0} \approx 376.73\;\Omega$. In the VCA framework, this possesses a literal mechanical identity: it is the \textit{acoustic impedance} of the discrete lattice, encoded into an electrical constant by the Topo-Kinematic mapping.

\subsection{Derivation from the Discrete LC Ladder}

A single lattice cell of pitch $\ell_{node}$ carries distributed inductance and capacitance per unit length equal to the vacuum permeability and permittivity:
\begin{resultbox}{Per-Cell Lumped Elements}
\begin{equation}
    L_{cell} = \mu_0\, \ell_{node}, \qquad C_{cell} = \epsilon_0\, \ell_{node}
    \label{eq:cell_elements}
\end{equation}
\end{resultbox}

Evaluating numerically from the physics engine constants ($\ell_{node} = 3.862 \times 10^{-13}$ m):
\begin{align}
    L_{cell} &= 1.257 \times 10^{-6} \times 3.862 \times 10^{-13} = 4.855 \times 10^{-19} \text{ H} \\
    C_{cell} &= 8.854 \times 10^{-12} \times 3.862 \times 10^{-13} = 3.419 \times 10^{-24} \text{ F}
\end{align}
The characteristic impedance of each lattice cell is:
\begin{resultbox}{Scale-Invariant Characteristic Impedance}
\begin{equation}
    Z_{cell} = \sqrt{\frac{L_{cell}}{C_{cell}}} = \sqrt{\frac{\mu_0\, \ell_{node}}{\epsilon_0\, \ell_{node}}} = \sqrt{\frac{\mu_0}{\epsilon_0}} \equiv Z_0 \approx 376.73\;\Omega
    \label{eq:z0_cell}
\end{equation}
\end{resultbox}

The lattice pitch cancels identically. This is the fundamental reason $Z_0$ is a universal constant: it is a property of the node-to-node impedance ratio of the lattice, independent of the absolute scale $\ell_{node}$. Every cell, at every location in the universe, presents the same $376.73\;\Omega$ characteristic impedance.

\subsection{Signal Propagation Velocity}

The group velocity of a signal through a lumped LC ladder is:
\begin{resultbox}{Propagation Velocity from Discrete Components}
\begin{equation}
    v_g = \frac{\ell_{node}}{\sqrt{L_{cell}\, C_{cell}}} = \frac{\ell_{node}}{\sqrt{\mu_0\, \epsilon_0\, \ell_{node}^2}} = \frac{1}{\sqrt{\mu_0\, \epsilon_0}} \equiv c
    \label{eq:c_from_lc}
\end{equation}
\end{resultbox}

The invariant speed of light is structurally identical to the slew rate of a discrete LC transmission line. No continuous medium is assumed; $c$ emerges from lumped elements.

\subsection{Mechanical Acoustic Impedance}

Applying the Topo-Kinematic identity (Section~\ref{sec:topo_kinematic}), the mechanical impedance equivalent is:
\begin{resultbox}{Mechanical Acoustic Impedance of the Vacuum}
\begin{equation}
    Z_{mech} = \xi_{topo}^{2}\, Z_0 = (4.149 \times 10^{-7})^2 \times 376.73 \approx 6.485 \times 10^{-11} \text{ kg/s}
    \label{eq:z_mech}
\end{equation}
\end{resultbox}

This value represents the absolute force-per-velocity ratio of a single lattice node. It is the fundamental quantum of mechanical impedance in the substrate.

\subsection{Impedance Across Physical Regimes}

Table~\ref{tab:impedance_regimes} summarises how $Z_0$ behaves across the three operating regimes of the vacuum, as established in Volume~I.

\begin{table}[H]
\centering
\caption{Characteristic impedance across the three vacuum operating regimes. Gravity preserves $Z_0$ (stealth); particle cores and event horizons destroy it ($Z \to 0$).}
\label{tab:impedance_regimes}
\small
\begin{tabularx}{\textwidth}{@{}@{@{}}l c c c c@{}}
\toprule
\textbf{Regime} & $\mu_{eff}$ & $\epsilon_{eff}$ & $Z = \sqrt{\mu_{eff}/\epsilon_{eff}}$ & $\Gamma$ \\
\midrule
Linear vacuum             & $\mu_0$                & $\epsilon_0$                & $Z_0 = 376.73\;\Omega$  & $0$ \\
Gravity well ($n \gg 1$)  & $n\, \mu_0$            & $n\, \epsilon_0$            & $Z_0 = 376.73\;\Omega$  & $0$ \\
Particle core ($\Delta\phi \to \alpha$) & $\to 0$ (Meissner) & $\to 0$ (dielectric collapse) & $\to 0\;\Omega$ & $-1$ \\
Event horizon             & $\to 0$ (saturated)    & $\to 0$ (saturated)         & $\to 0\;\Omega$          & $-1$ \\
\bottomrule
\end{tabularx}
\end{table}

\noindent The critical distinction: gravity scales $\mu$ and $\epsilon$ \textit{symmetrically} ($n \times n$), preserving $Z_0$ and producing zero reflection. Topological saturation (particles, event horizons) drives both to zero \textit{asymmetrically} via Axiom~4, collapsing $Z$ and creating perfect mirrors ($\Gamma = -1$). This is why gravity wells are RF-transparent (stealth) while particle boundaries and event horizons are perfect reflectors.\section{Gravitational Stealth (S-Parameter Analysis)}

In classical RF engineering, when a wave transitions into a denser physical medium, the refractive index ($n$) rises asymmetrically, forcing the characteristic impedance to drop. This impedance mismatch causes the signal to partially reflect, measured logarithmically as Return Loss ($S_{11}$). This introduces a profound paradox for analog gravity models: \textit{If a gravity well represents a physical increase in the localized optical density of the vacuum, why does light seamlessly enter a black hole without scattering or reflecting off the boundary?}

In the VCA transmission line model, macroscopic gravity operates strictly as a 3D Volumetric Compression of the Chiral LC Network \cite{misner1973}. This localized geometric crowding proportionately and \textit{symmetrically} increases both the effective inductive mass density ($\mu_{local} = n(r) \cdot \mu_0$) and the capacitive compliance ($\epsilon_{local} = n(r) \cdot \epsilon_0$). Evaluating the Characteristic Impedance of the vacuum down to the extreme metric divergence of an Event Horizon ($r \to R_s$) reveals a perfect mathematical invariant:
\begin{equation}
Z_{local}(r) = \sqrt{\frac{\mu_{local}}{\epsilon_{local}}} = \sqrt{\frac{n(r) \cdot \mu_0}{n(r) \cdot \epsilon_0}} = \sqrt{\frac{\mu_0}{\epsilon_0}} \equiv Z_0 \approx 376.73\ \Omega
\end{equation}

The substrate is mathematically and perfectly Impedance-Matched to itself everywhere, absolutely regardless of extreme gravitational strain. Because the spatial derivative of the impedance remains strictly zero ($\partial_r Z_0 = 0$), the Reflection Coefficient ($\Gamma$) is mathematically forced to zero. The universe structurally possesses an \textbf{$S_{11}$ Return Loss of $-\infty$ dB}. This provides the exact continuum-mechanics mechanism for why localized gravitational gradients act as perfect RF-absorbing stealth structures rather than optical mirrors.

\medskip\noindent\textbf{The branch-vertex counterpoint --- an intrinsically mismatched reciprocal 3-port.} The symmetric \emph{volumetric} scaling above matches a \emph{2-port} boundary. A \emph{3-way branch} does not have that luxury. At the chiral \emph{srs} vacuum-net vertex (coordination $z=3$), a wave down one bond sees the other two in parallel ($Z_0/2$), so the bare junction reflects $\Gamma = (2-z)/z = -\tfrac13$ ($|\Gamma|^2 = \tfrac19$, about $-9.5$~dB) --- a pure counting fact (one bond feeding two), reactively back-scattering (Axiom~3, not dissipation). The floor $|S_{11}| \ge \tfrac13$ is the classic matched-lossless-reciprocal-3-port theorem (Pozar-class), confirmed at the vertex: \emph{no} lossless reciprocal bore reaches the $S_{11} \to -\infty$~dB match a symmetric 2-port enjoys. The \emph{only} matched 3-port is non-reciprocal --- the ideal circulator ($S_{11}=0$, unitary, $C_3$) --- which requires a time-reversal-breaking bias ($I4_132$ chirality is only parity-breaking); whether the vacuum supplies one is PENDING-GRANT (asserted nowhere). So the gravity well is RF-transparent and the saturated particle core is a perfect $\Gamma=-1$ mirror, but a bare 3-way vacuum junction sits between them: a hard $\tfrac13$ reciprocal reflection floor unless a circulator bias intervenes. Canonical at \kbleaf{ave-kb/vol1/operators-and-regimes/ch6-universal-operators/srs-vertex-scattering.md} (\texttt{clm-v3port}, \texttt{clm-bore2x}; X38/X37).

\begin{figure}[htbp]
    \centering
    \includegraphics[width=\textwidth]{log_impedance_s_parameters.png}
\caption{\textbf{S-Parameter Analysis of a Gravity Well.} Top: As a wave approaches a gravitational core, the density $n(r)$ diverges. Because analog macroscopic gravity compresses volumetric space, it scales $L$ and $C$ symmetrically, ensuring the Characteristic Impedance ($Z_0$) remains perfectly invariant. Bottom: If gravity behaved like an unmatched optical dielectric, the resulting impedance drop would generate massive reflection ($S_{11} > -10$ dB). The symmetric volumetric scaling of the AVE EFT forces $S_{11} \to -\infty$ dB, providing the precise mechanism for why intense gravity wells do not act as RF mirrors.}
    \label{fig:s_parameters}
\end{figure}

\subsection{The Substrate Transmission Line (Emergence of \texorpdfstring{$c$}{c})}
To computationally prove that macroscopic Special Relativity emerges deterministically from these discrete components, the 1D spatial vacuum grid as a cascaded LC transmission line. By normalizing the discrete Inductors ($\mu_0 \ell_{node}$) and Capacitors ($\epsilon_0 \ell_{node}$) to the hardware pitch, the injection of a transient topological voltage pulse confirms that the signal propagates through the discrete components at exactly the continuous group velocity $v_g = 1/\sqrt{LC} \equiv c$. The continuous, invariant speed of light is mathematically identically the macroscopic slew-rate of a discrete transmission line.

\begin{figure}[htbp]
    \centering
    \includegraphics[width=\textwidth]{condensate_transmission_line.png}
\caption{\textbf{The EFT Transmission Line.} A time-domain simulation of a discrete 100-node vacuum grid. By cascading the discrete inductive mass and capacitive compliance of the analog lattice, the signal propagates flawlessly at $v_g = c$, proving that continuous spacetime kinematics emerge natively from lumped-element circuit analysis.}
    \label{fig:transmission_line}
\end{figure}

\subsection{The Horizon Mirror: Predicting Black Hole Echoes}

While the bulk continuous gravity well remains perfectly impedance-matched ($Z = Z_0$), the exact mathematical boundary of the Event Horizon represents a profound physical discontinuity. 

As established in Chapter 9, the Event Horizon is strictly defined as the radius where the volumetric tensor strain reaches the absolute Axiom 4 dielectric saturation limit ($\Delta\phi \to \alpha$). At this precise topological boundary, the effective capacitance of the macroscopic metric diverges to infinity ($C \to \infty$). 

Consequently, the characteristic impedance of the spacetime metric exactly at the event horizon mathematically collapses to zero ($Z_{EH} \to 0\,\Omega$). Evaluating the reflection coefficient between the deep gravity well ($376.7\,\Omega$) and the event horizon ($0\,\Omega$) yields:

\begin{equation}
\Gamma_{EH} = \frac{Z_{EH} - Z_{0}}{Z_{EH} + Z_{0}} = \frac{0 - 376.7}{0 + 376.7} = -1
\end{equation}

This reveals that while a gravity well is "stealthy" to approaching waves, the Event Horizon itself acts as a macroscopic, perfect topological mirror. Infalling energy that reaches the absolute saturation limit undergoes a perfect $180^\circ$ phase inversion and reflects outward. This explicitly predicts the existence of \textbf{Black Hole Echoes}—post-merger gravitational wave reflections currently hypothesized by advanced quantum gravity models—providing a strict, testable falsification metric for the AVE framework via future LIGO/LISA observations.

\section{The Periodic Table: Topological SPICE Mappings}
To decisively bridge the theoretical kinematic behavior of a single LC knot into complex, multi-body engineering, the AVE framework mathematically derives the entire structural sequence of atomic elements. 

As fully derived in the supplementary \textit{Periodic Table of Knots}, the 3D stable geometry of an atomic nucleus is deterministically bounded by these exact inductive LC matrices ($M_{ij}$). By mathematically mapping protons ($6^3_2$ topological lattice defaults) into continuous Spatial SPICE networks, the engine computes exactly how atomic geometry collapses into absolute symmetric extremes (the noble gases and base metallics) versus partial, highly reactive asymmetric valences (the halogens and alkali metals).

% 🔴 METHODOLOGY DISCLOSURE 2026-08-02 (Rule 12 KEEP-BOTH; the Silicon-28 paragraph and the
%   "strictest possible verification" sentence below are PRESERVED VERBATIM -- no number is
%   edited and no number is replaced). Ruled authority: ave-kb/vol1/operators-and-regimes/
%   ch5-universal-spatial-tension/scale-invariance.md:35 -- "Methodology disclosure (per Vol 6
%   introduction). Within the (Z,A)-forced bipyramid topology, the inter-alpha distance
%   R_bipyramid is the single fitted scalar adjusted per nucleus so the pairwise summation
%   reproduces the CODATA mass" -- and :41 -- "The reported <0.001% residual is the optimizer's
%   convergence tolerance on R, not an independent prediction error on the mass." Corroborated
%   by ave-kb/vol1/claim-quality.md:531 ("NOT zero-parameter ab initio") and
%   ave-kb/vol6/period-2/neon/index.md:21 ("structurally zero by construction").
%   FLAG-DON'T-FIX (routed, NOT adjudicated here; corrected 2026-08-02 per #827 audit F1/F5/F6):
%   the printed R_bipyramid = 80.174d agrees with the engine (src/scripts/vol_6_periodic_table/
%   simulations/simulate_element.py:390, r_bipyr = 80.174370*d) but diverges from the KB's Si-28
%   row at ave-kb/vol6/appendix/geometric-inevitability/platonic-progression.md:24 (R ($d$) =
%   83.0) -- same quantity, same units; the KB row is the staleness candidate (routed). That
%   leaf's 0.0002% column is a geometric residual to assumed Platonic coordinates (its :14 scope
%   correction), NOT a mass-mapping error -- do not compare it to the 0.0000% below. [A prior
%   limb here asserted the printed "378 coupled inductor nodes" diverges from the KB's 21
%   inter-alpha pairs; WITHDRAWN -- the KB carries both, reconciled (claim-quality.md:629 "378
%   SPICE coupled inductors"; period-3/silicon/ee-equivalent.md:10 "378-parameter SPICE matrix
%   decomposes into 21 inter-alpha links"). Residual wording nit only: print says "inductor
%   nodes", KB says "coupled inductors".]
%
%   ---- CORRECTION 2026-08-03 (Rule 12 quote-and-date; the block above is preserved
%   verbatim, unedited). The sentence above reading
%       "the KB row is the staleness candidate (routed)"
%   is STRUCK. It elects a side, and that election CONTRADICTS the ratified rev-2
%   correction fragment for this exact item -- _orchestration/docket-entries/
%   2026-08-02-mr-no-kb-home-correction.md:5-6. Three receipts, all re-verified by
%   content 2026-08-03:
%     (1) BOTH VALUES ARE FITS, neither is "the derived authority."
%         simulate_element.py:10-15 (its own scope note): the radii "were 1-parameter
%         Nelder-Mead fits against CODATA mass targets"; :388-389 "Numerically optimized
%         to ~80.174370d ... matching the 26053.188074 MeV empirical nuclear target".
%         solve_silicon.py:5-9: "a 1-parameter inverse-problem solve, NOT a forward
%         prediction". symmetric-core-collapse.md:10 fits 83.0d to the SAME 26053.188
%         MeV target. Engineering-choice fit tags apply on both sides; "the engine
%         computed it, therefore it is canonical" is not an available argument.
%     (2) THERE ARE TWO MODELS, not one solver with a stale mirror.
%         ave-kb/vol6/framework/computational-mass-defect/semiconductor-nuclear-analysis.md
%         :20 names simulate_element.py the "Bare K/r Model"; :22 names
%         semiconductor_binding_engine.py the "Semiconductor Junction Model"; :26-32
%         reconciles them per element (Si-28 row at :32 -- bare 85.6d vs semiconductor
%         83.0d); :34 rules "The semiconductor model R values are the definitive
%         mass-validated quantities."
%     (3) platonic-progression.md:24's 83.0 is therefore the SEMICONDUCTOR /
%         DEFINITIVE value -- it MATCHES the ruled-definitive column and is NOT stale.
%   THE ACTUAL DIVERGENCE PAIR is a different pair entirely: the KB BARE cell
%   (85.6d, semiconductor-nuclear-analysis.md:32) vs the engine bare hardcode
%   (80.174370d, simulate_element.py:390) -- same model, same quantity, same units,
%   two values. And the print sentence below quotes a BARE-model R (80.174d) under a
%   regime where the KB rules the SEMICONDUCTOR R definitive.
%   WHAT IS ROUTED (open, NOT adjudicated here): the METER CHOICE -- is the print
%   quoting the wrong model for this passage (bare where semiconductor is ruled
%   definitive), or has the KB bare cell gone stale against the engine? That is a
%   Grant-level ruling and neither this comment nor this commit takes it. No printed
%   number is edited here; only the wrong election is struck.
%   Dependent set if any value moves (from the fragment :9): 83.0d sites --
%   platonic-progression.md:24, semiconductor-nuclear-analysis.md:32,
%   vol6/period-3/index.md:22, vol6/period-3/silicon/index.md:18,:31,
%   vol6/period-3/silicon/symmetric-core-collapse.md:10; bare side --
%   simulate_element.py:388-390, and the print sentence below. ----
\noindent\textbf{Methodology disclosure (2026-08-02, per the Vol~6 fit/predict split; the paragraph below is preserved per Rule~12 KEEP-BOTH).} The Silicon-28 result below reports a residual as though it were a prediction error. Per the ruled disclosure, within the $(Z,A)$-forced bipyramid topology the inter-alpha distance $R_{bipyramid}$ is \textbf{the single fitted scalar adjusted per nucleus} so that the pairwise summation reproduces the CODATA mass. What is \emph{predicted} (axiom-derived) is the cluster topology as a function of $(Z,A)$ via minimum-impedance packing, the coupling $K_{mutual} = m_p c^2$, and the parameter count (one scalar per nucleus); what is \emph{fitted} is the numerical value of $R_{bipyramid}$. Consequently the quoted mass-mapping residual is \textbf{the optimizer's convergence tolerance on $R$, not an independent prediction error on the mass} --- an exactly-zero residual on a closed shell is structurally zero by construction, and the pipeline is \textbf{not} zero-parameter \emph{ab initio}. The falsifiable structural content is the topology and the one-scalar-per-nucleus parameter count (against the $\sim 5$ parameters per nucleus of liquid-drop or shell-model fits), not the residual. Read the ``strictest possible verification'' sentence below in that scope. Canonical: \kbleaf{ave-kb/vol1/operators-and-regimes/ch5-universal-spatial-tension/scale-invariance.md}.

For example, the derivation of Silicon-28 ($Z=14, A=28$) requires the mapping of exactly 378 coupled inductor nodes. The solver algebraically targets the empirical CODATA nuclear mass~\cite{codata2018}, forcing the geometry to converge. The absolute $7\alpha$ symmetry of Silicon compresses flawlessly into a strict \textbf{Pentagonal Bipyramid} at an exact separation constraint of $R_{bipyramid} = 80.174d$, yielding a \textbf{$0.0000\%$ mass mapping error}. 

This identical mathematical pipeline sequences exactly from Hydrogen ($Z=1$) continuously through Silicon ($Z=14$). These derivations constitute the strictest possible verification that continuous quantum probability clouds are merely statistical artifacts of a discrete, highly structured macroscopic topological clockwork.

\subsection*{Impedance Boundary Regime Classification}
\begin{center}
\small
\renewcommand{\arraystretch}{1.2}
\begin{tabularx}{\textwidth}{@{}lccX@{}}
\toprule
\textbf{Region} & \textbf{$Z_{local}$} & \textbf{$\Gamma$} & \textbf{Physical Character} \\
\midrule
Free space             & $376.7\,\Omega$ & $0$   & Regime I: lossless propagation \\
Gravity well ($r > r_s$) & $376.7\,\Omega$ (invariant) & $0$   & Regime I: stealth (symmetric $\epsilon\mu$ scaling) \\
Event horizon ($r = r_s$) & $0\,\Omega$  & $-1$ & Regime III: perfect mirror \\
Interior ($r < r_s$)     & Undefined    & ---   & Pre-geometric plasma \\
\bottomrule
\end{tabularx}
\end{center}\section{Topological Defects as Resonant LC Solitons}

As established in prior chapters, a fundamental particle is a stable topological defect---a highly
tensioned phase vortex permanently locked into the discrete graph structure. In classical
electrical engineering, a localized, trapped electromagnetic standing wave that permanently
cycles reactive energy without radiative loss is defined as a \textbf{Resonant LC Tank Circuit}.

By applying the Topo-Kinematic mapping to the electron's rest mass, its equivalent localized
Inductance evaluates to $L_e \equiv \xi_{topo}^{-2} m_e$. The local lattice compliance acts as the restoring
capacitor ($C_e \equiv \xi_{topo}^2 k^{-1}$).

\subsection{Recovering the Virial Theorem and $E=mc^2$}

This mapping is rigorously verified this structural mapping by evaluating the stored energy of the
resonant soliton. In an ideal LC tank, the peak internal dynamic (inductive) energy is defined
as $E_{mag} = \frac{1}{2} L_e I_{max}^2$. Substituting the hardware velocity limit ($I_{max} = \xi_{topo} c$) evaluates to:

\begin{equation}
E_{mag} = \frac{1}{2} (\xi_{topo}^{-2} m_e)(\xi_{topo} c)^2 = \frac{1}{2} m_e c^2
\end{equation}

In a stable LC resonant soliton, the classical Virial Theorem rigidly dictates that the
capacitive (electric/strain) energy stored in the static topological twist of the core must exactly
equal the inductive kinetic energy ($E_{elec} = E_{mag} = \frac{1}{2}m_e c^2$). Summing the two isolated energy
ledgers perfectly recovers $E_{total} = m_e c^2$ \cite{einstein1916}. Einstein's mass-energy equivalence principle is
mechanically and mathematically identical to the Total Stored Electrical Energy of a classical
macroscopic Resonant LC Tank Circuit ringing natively within the analog vacuum metric.

\subsection{Total Internal Reflection: The Confinement Bubble}

A fundamental requirement for any discrete particle (soliton) model is explaining why the
localized wave-packet does not instantly disperse its stored energy into the ambient vacuum.
In the AVE framework, this geometric stability is mathematically guaranteed by the extreme
flux crowding at the particle's boundary, which generates a perfect macroscopic impedance
mismatch.

Unlike the symmetric volumetric compression of macroscopic gravity (which keeps $Z_{0}$
perfectly invariant, preventing scattering), the localized topological twist of a particle core
induces extreme dielectric saturation. As the local topological strain ($\Delta\phi$) approaches the
Axiom 4 hardware limit ($\alpha$), the effective geometric capacitance (compliance) of the boundary
nodes diverges to infinity:

\begin{equation}
\lim_{\Delta\phi \to \alpha} C_{eff}(\Delta\phi) = \lim_{\Delta\phi \to \alpha} \frac{C_{0}}{\sqrt{1-\left(\frac{\Delta\phi}{\alpha}\right)^{2}}} = \infty
\end{equation}

Because the characteristic impedance of a spatial cell is dictated by $Z=\sqrt{L/C}$, this
massive spike in boundary capacitance drives the localized impedance of the particle boundary
strictly to zero:

\begin{equation}
\lim_{C_{eff} \to \infty} Z_{core} = \lim_{C_{eff} \to \infty} \sqrt{\frac{\mu_{0}}{C_{eff}}} = 0\,\Omega
\end{equation}

In standard wave mechanics, the Reflection Coefficient ($\Gamma$) governing the transmission of
energy across a boundary is defined by the impedance differential between the two media.
Evaluating the boundary between the saturated particle core ($0\,\Omega$) and the unperturbed
ambient vacuum ($Z_{0}\approx376.7\,\Omega$) yields:

\begin{equation}
\Gamma=\frac{Z_{core}-Z_{0}}{Z_{core}+Z_{0}}=\frac{0-376.7}{0+376.7}=-1
\end{equation}

A reflection coefficient of $\Gamma=-1$ constitutes a \textbf{Perfect Short-Circuit Boundary}.

This mathematical limit proves that 100\% of the kinetic energy attempting to radiate
outward from the saturated flux tube hits this impedance wall, undergoes a perfect $180^{\circ}$ phase
inversion, and reflects internally. Mechanically, the nodes at the saturation boundary are
geometrically jammed at the absolute hard-sphere exclusion limit. The local phase velocity
($c_{local}=1/\sqrt{LC}$) strictly collapses to zero, creating a hyper-rigid, localized envelope. The
particle dynamically weaves its own perfect topological mirror, forming an impenetrable,
hyper-highly-reluctant ``Local Bubble'' that perfectly confines the internal LC resonance without
radiative loss.

\begin{figure}[h]
    \centering
    \includegraphics[width=0.62\textwidth]{electron_lattice_net_S.png}
    \caption{\textbf{The ``Local Bubble'' on the real vacuum lattice (engine-exact).} The seeded $(2,3)$ winding on the chiral \emph{srs} net (degree-$3$, $I4_132$ / (10,3)-a; $L=6$, $1728$ interior $z=3$ nodes), projected to the $(x',y')$ plane. Node colour is the Axiom-4 saturation kernel $S(A)=(1-A^2)^{0.5}$ at the winding amplitude field; the ring of open markers is the set of $19$ \textbf{wall nodes at yield} ($A>0.9$, $S\!\to\!0.045$) --- the $Z_{core}\to0$, $\Gamma=-1$ perfect short-circuit shell derived above, rendered on the canonical carrier. The teal cycle is the $\Delta b_1=+1$ meridian doorway (linking $\pm1$ with the winding core, verified). Axes in $a_{\mathrm{cell}}$; the in-plane projection is presentation-layer, all node/bond/winding/kernel values are engine-exact. Driver: \texttt{src/scripts/viz/electron\_lattice\_scene.py} (\texttt{ave.core.chiral\_lattice.build\_srs\_net} + \texttt{srs\_cage\_winding} + \texttt{graded\_vacuum\_network.saturation\_kernel}); provenance ledger \texttt{viz/README.md} Visual~1.}
    \label{fig:electron_lattice_local_bubble}
\end{figure}

\textbf{Deriving the QCD Linear Potential:} Furthermore, this provides the strict deter-
ministic mechanism for Strong Force flux collimation. Rather than radiating isotropically
($1/r^{2}$), the energy traveling between nucleons undergoes Total Internal Reflection (TIR) off
the impedance walls of the highly strained vacuum, acting as a Topological Fiber-Optic Cable.

By applying Gauss's Law to a confined 1D cylinder of constant cross-sectional area,
the electric flux density ($D$) mathematically cannot spread radially outward. The electric
flux remains perfectly constant along the entire length of the tube, absolutely regardless
of separation distance. Consequently, the restorative force ($F(r) = \text{constant}$) inherently
generates the exact \textbf{Linear Confinement Potential} ($V(r)\propto r$) empirically observed in
Quantum Chromodynamics. The phenomenological ``MIT Bag Model'' is directly exposed
as a macroscopic impedance wall woven natively by the non-linear varactor limits of the
continuous vacuum.

\subsection{The Mechanical Origin of the Pauli Exclusion Principle}

The establishment of the saturated particle boundary as a perfect topological mirror ($\Gamma = -1$)
provides a rigorous, continuous-mechanical derivation for the Pauli Exclusion Principle. 

In standard quantum mechanics, the inability of fermions to occupy the same quantum
state is treated as an abstract statistical postulate. In the AVE framework, it is an unavoidable
consequence of classical macroscopic impedance boundaries. 

When massless Bosons (photons) propagate, they act as linear transverse shear waves.
Because they do not possess a static inductive core, they do not geometrically saturate
the dielectric lattice ($\Delta\phi \ll \alpha$). The local metric impedance remains perfectly matched at
$Z_{0} \approx 376.7\,\Omega$. With a reflection coefficient of $\Gamma \approx 0$, boson waves pass cleanly through one
another, permitting infinite superposition.

Conversely, Fermions are massive topological defects bounded by strictly saturated $Z_{core} =
0\,\Omega$ envelopes. If two fermions are forced into the same spatial volume, their boundaries
collide. Because both boundaries possess a reflection coefficient of strictly $\Gamma = -1$, their
internal localized wave-functions cannot mathematically penetrate one another. The kinetic
energy of Fermion A perfectly reflects off the infinite-compliance wall of Fermion B. The Pauli
Exclusion Principle is therefore physically identical to the hard-sphere collision of perfectly
impedance-mismatched dielectric bubbles.

\section{Real vs. Reactive Power: The Orbital Friction Paradox}

A historical and persistent critique of analog inductive spacetime models is the ``Friction Paradox'':
\textit{If a planet is physically moving through a dense spatial substrate, why doesn't inductive drag
drain its kinetic energy, causing its orbit to decay over cosmological timescales?}

Within the VCA framework, this paradox is resolved flawlessly by rigorously distinguishing
between non-conservative inductive drag and conservative AC Power Analysis. As established
in Chapter 11, exceeding the Dielectric Saturation limit ($\tau > \tau_{yield}$) does not merely result in a
classical highly-reluctant network; it triggers an avalanche dielectric phase-transition. The local metric
structurally melts into an irrotational, continuous quantum network. Because this continuous
melted phase mathematically cannot support transverse shear vectors, the localized inductive
mutual inductance strictly collapses to zero ($\eta \to 0$). Therefore, the anti-parallel inductive drag force
($F_{drag}$) mathematically evaluates to exactly zero Newtons \cite{flyby2008}.

With non-conservative drag structurally eliminated, evaluating the remaining thermody-
namic interaction using electrical engineering power principles. Total apparent power ($S$) is
divided into two distinct components depending on the phase angle ($\theta$) between Voltage ($V$)
and Current ($I$):
\begin{enumerate}
\item \textbf{Real Power ($P$):} Measured in Watts. $P = VI \cos(\theta)$. This represents energy physically
dissipated from the system.
\item \textbf{Reactive Power ($Q$):} Measured in Volt-Amperes Reactive (VARs). $Q = VI \sin(\theta)$.
This represents energy conservatively exchanged back and forth without permanent
dissipation.
\end{enumerate}

By applying the Topo-Kinematic Identity to the remaining conservative interactions, the
radial Gravitational Force vector acts identically as the AC Voltage ($V_{substrate} \propto F_g$), and
the tangential Orbital Velocity vector acts as the AC Current ($I_{substrate} \propto v_{orb}$). In a stable,
circular planetary orbit, the radial gravitational force vector is perfectly and mathematically
orthogonal ($90^\circ$) to the tangential velocity vector. Therefore, the phase angle between the
vacuum Voltage and Current is exactly $\theta = 90^\circ$. 

Evaluating the Real Power physically dissipated by the planetary body into the vacuum
network via the conservative gravity well yields:
\begin{equation}
P_{real} = F_g \cdot v_{orb} \cdot \cos(90^\circ) \equiv 0\ \text{Watts}
\end{equation}

Because inductive drag is neutralized by the dielectric phase transition, and the remaining
gravitational coupling is purely orthogonal, the orbiting body experiences absolutely zero
macroscopic energy dissipation. A stable planetary orbit is the macroscopic mechanical
equivalent of a \textbf{Lossless LC Tank Circuit} operating purely in the reactive power domain.

\subsection*{Power Domain Classification}
\begin{center}
\small
\renewcommand{\arraystretch}{1.2}
\begin{tabularx}{\textwidth}{@{}lcccX@{}}
\toprule
\textbf{System} & \textbf{$\theta$} & \textbf{$P_{real}$} & \textbf{$Q_{reactive}$} & \textbf{AVE Interpretation} \\
\midrule
Stable orbit     & $90^\circ$ & $0$ W & $F_g \cdot v_{orb}$ & Lossless LC tank \\
Inspiral (GW)    & $\neq 90^\circ$ & $> 0$ & Decreasing & Radiative damping \\
Electron orbital & $90^\circ$ & $0$ W & $m_e c^2 \cdot \alpha$ & Quantized reactive shell \\
Photon propagation & $0^\circ$ & $P_{rad}$ & $0$ & Pure real (travelling wave) \\
\bottomrule
\end{tabularx}
\end{center}\section{Substrate IMD Spectroscopy: The Harmonic Fingerprint}
\label{sec:imd}

% 🔴 RECONCILED 2026-06-21 (Rule 12): the IMD drive table, the "~30\% / ~13 kV / tabletop" framing, and the "QED predicts a sextic ($V^6$) scaling" discriminator below are PRESERVED verbatim; the scope-correction note recasts their framing per the per-node-vs-apparatus discipline (KB sibling vacuum-impedance-mirror.md, 2026-06-04) and the COEFFICIENT discriminator (KB clm-pp3qwf, commit ad26d357). Lockstep with the KB twin ave-kb/vol4/circuit-theory/ch1-vacuum-circuit-analysis/intermodulation-distortion.md; sibling to the 2026-06-20 birefringence walk-back in Ch.~11.
\noindent\textbf{Scope correction (2026-06-21, per-node-vs-apparatus + IM3 discriminator reconciliation; body below preserved per Rule 12).} Three framings in this section are corrected. (1) The drive-level table reads the \textbf{apparatus} drive voltage (kV across a $\sim$100 µm gap) as if it were the \textbf{per-node} strain $A = E_{local}/E_{yield}$ (with $E_{yield} = V_{yield}/\ell_{node} \approx 1.13\times10^{17}$ V/m, a per-node \emph{field}); this overstates the Ax-4 kernel argument $A$ by $\sim 2.6\times10^8\times$ at a 100 µm gap. At the "$0.30$ / $13.10$ kV" row, the honest per-node strain is $A \approx \mathbf{1.2\times10^{-9}}$ (not $0.30$), giving $\text{IM3} \approx \mathbf{-360}$ dBc (not $-70$). (2) The "becomes measurable above $\sim 30\%$ of $V_{yield}$ ($\sim 13$ kV)" / "tabletop" foothold is \textbf{retracted}: on the correct per-node kernel the $-80$ dBc detection field is $E \approx 1.3\times10^{15}$ V/m --- \textbf{facility-class}, the same regime as the vacuum-birefringence test (Ch.~12, \S``The Vacuum Birefringence Limit''), not a lower-field benchtop foothold. (3) The "QED predicts a sextic ($V^6$) scaling" discriminator is \textbf{retracted} (KEEP-BOTH): both AVE and QED IM3 are \textbf{cubic-in-drive} ($\propto V^3$, both from a quartic $E^4$ Lagrangian $\to \chi^{(3)}$); QED's "$^6$" is the \emph{frequency} exponent $(\omega/m_ec^2)^6$ of the Euler-Heisenberg cross-section, not a voltage slope, so the IM3-vs-drive exponent reads $3$ for both and cannot discriminate. The real discriminator is the \textbf{coefficient} (an echo, $E_{yield} = \sqrt{\alpha}\,E_{crit}$), per KB \texttt{clm-pp3qwf} and the B1 row in \texttt{ave-kb/common/divergence-test-substrate-map.md}.

By modelling the universe as a non-linear dielectric network, the AVE framework makes a specific, falsifiable prediction absent from standard Quantum Electrodynamics: the vacuum should produce measurable \textbf{Intermodulation Distortion (IMD)} products when driven by sufficiently intense electromagnetic fields. This section derives the expected signature analytically and specifies the experimental parameters for detection.

\subsection{The Non-Linear Source Term}

\noindent\textbf{Sector re-keying correction (2026-08-03, $V_{yield} \to V_{snap}$; \texttt{def-vyvsn1} grade-fork RESOLVED $=$ T2, Grant 2026-06-30; the rest of this section is preserved per Rule~12 KEEP-BOTH).} This section keys the \emph{diverging} compliance $C_{vac} = C_0/S$ on $V_{yield} = \sqrt{\alpha}\,V_{snap} \approx 43.65$~kV. That keying is \textbf{corrected} --- this is the un-propagated remainder of the same ratified ruling already applied at \S``The Vacuum Varactor'' above, whose note is explicitly scoped \emph{``this note covers this subsection only.''} \textbf{(1) Which reactance.} The diverging $C_0/S$ is the \textbf{longitudinal-A1 bond compliance} (the stretch-reactance $1/k_a$), \emph{not} the transverse dielectric capacitance; the transverse permittivity $\varepsilon_{eff} = \varepsilon_0 S$ and the across-gap cell capacitance an LCR meter reads, $C_{diel} \propto S$, \emph{roll off} instead. A1 $\perp$ T2 are orthogonal reactances sharing the EE name ``capacitance''; identifying them is the genesis-24 double-count. \textbf{(2) Which voltage.} $V_{yield}$ is the \textbf{transverse Cosserat ($T_2$) self-trap wall}, not the A1 compliance bound; the A1 compliance diverges at the higher $V_{snap} = m_ec^2/e \approx 511$~kV, a factor $1/\sqrt{\alpha} = 11.706$ above $V_{yield}$. \textbf{(3) How to read this section.} Eq.~\ref{eq:varactor_imd} and the IP3 result (Eq.~\ref{eq:ip3}) are \textbf{updated in place}; the Taylor expansion below, Table~\ref{tab:imd_levels} and the falsification criterion keep their $V_{yield}$ token \textbf{preserved per Rule~12} and are read with $V_{yield} \to V_{snap}$. The dimensionless \emph{ratios} ($C_{vac}/C_0$, $S$, and the $V/V_{\cdot}$ column) are \textbf{unchanged}; only the absolute voltage column rescales by $1/\sqrt{\alpha}$, so the printed $43.65$~kV endpoint is the $V/V_{snap} = \sqrt{\alpha} \approx 0.085$ \emph{first} row of the corrected table, not its divergence. \textbf{(4) What this does NOT do.} It does \emph{not} revive the retracted tabletop foothold --- the 2026-06-21 scope correction above stands untouched on the \emph{per-node} argument, which is orthogonal to which apparatus voltage normalizes the ratio (honest detection field remains facility-class $E \approx 1.3\times10^{15}$ V/m) --- and it does not touch the KEEP-BOTH cubic-vs-sextic retraction. Canonical: \kbleaf{ave-kb/vol4/circuit-theory/ch1-vacuum-circuit-analysis/nonlinear-vacuum-capacitance.md} (\texttt{def-vyvsn1}):18, \kbleaf{ave-kb/vol4/falsification/ch12-falsifiable-predictions/ee-bench-plateau.md}:18--20; KB twin re-keyed in lockstep at \kbleaf{ave-kb/vol4/circuit-theory/ch1-vacuum-circuit-analysis/intermodulation-distortion.md}.

Standard QED models the vacuum as a linear medium at low energies, predicting that photon-photon scattering occurs only via extraordinarily weak perturbative quantum fluctuations (Euler-Heisenberg, $\sigma \propto \alpha^4$). The AVE framework replaces this with a deterministic, classical non-linearity: the squared geometric saturation limit (Axiom~4) imposes a macroscopic varactor on the dielectric constant of the vacuum:
\begin{resultbox}{Vacuum Varactor (Axiom~4)}
% [RE-KEYED 2026-08-03 per the sector-re-keying correction above -- prior form
%  preserved per Rule 12]: this equation read
%    "C_{vac}(V) = \frac{C_0}{\sqrt{1 - (V/V_{yield})^2}}, \qquad
%     V_{yield} \approx 43.65 \text{ kV}"
%  Token swapped V_yield -> V_snap; the divergent C_0/S is the longitudinal-A1
%  compliance, bounded at V_SNAP = 510,998.95 V (src/ave/core/constants.py),
%  not at the transverse-T2 wall V_YIELD = 43,651.85 V. Functional form
%  unchanged.
\begin{equation}
    C_{vac}(V) = \frac{C_0}{\sqrt{1 - (V/V_{snap})^2}}, \qquad V_{snap} \approx 511 \text{ kV}
    \label{eq:varactor_imd}
\end{equation}
\end{resultbox}

\begin{figure}[h]
    \centering
    \includegraphics[width=0.9\textwidth]{vacuum_dielectric_saturation.png}
\caption{\textbf{The Squared Dielectric Saturation Limit.} Unlike standard perturbative QED (Dashed Black), the AVE substrate (Red) imposes a macroscopic geometric varactor asymptote at $V_{yield}$. This classical structural non-linearity is the specific source of the macroscopic intermodulation products predicted by the AVE framework.}
    \label{fig:dielectric_saturation}
\end{figure}

\subsection{Third-Order IMD from Taylor Expansion}

The Taylor expansion of the varactor about $V = 0$ (Eq.~\ref{eq:varactor_taylor}) yields the non-linear susceptibility:
\begin{equation}
    C_{vac}(V) \approx C_0 \left[1 + \underbrace{\frac{1}{2}\left(\frac{V}{V_{yield}}\right)^{\!2}}_{\text{2nd order}} + \underbrace{\frac{3}{8}\left(\frac{V}{V_{yield}}\right)^{\!4}}_{\text{4th order}} + \cdots\right]
\end{equation}

When driven by a dual-tone signal $V(t) = V_1 \cos(\omega_1 t) + V_2 \cos(\omega_2 t)$, the squared term generates 2nd-order products at $\omega_1 \pm \omega_2$, and the quartic term generates the critical \textbf{3rd-order intermodulation} (IM3) products:
\begin{resultbox}{3rd-Order Intermodulation Products}
\begin{equation}
    f_{IM3} = 2f_1 - f_2 \quad \text{and} \quad 2f_2 - f_1
    \label{eq:im3_frequencies}
\end{equation}
\end{resultbox}

These IM3 tones fall \textit{close} to the original drive frequencies (unlike harmonic products at $2f$, $3f$ which fall far out of band), making them the most experimentally accessible non-linear signature.

\medskip\noindent\emph{Parity cross-ref (\texttt{clm-invmtr}, body preserved).} The 2nd-order products at $\omega_1 \pm \omega_2$ named above live in the \textbf{capacitance envelope} $C(t)$ (even in $V$), which is \emph{not} the radiated line: sub-yield, the difference tone $\omega_1 - \omega_2$ is inversion-symmetry-\textbf{forbidden} in the radiated observable (the odd-force parity theorem introduced with the metric varactor above; \kbleaf{ave-kb/common/universal-saturation-kernel-catalog.md}, \texttt{clm-invmtr}). The IM3 line $2f_1 - f_2$ ($m+n$ odd) is parity-\textbf{allowed} and unaffected.

\subsection{Predicted IM3 Amplitude}

For a standard varactor with $C(V) = C_0 (1 - V/V_{br})^{-1/2}$, the 3rd-order intercept point (IP3) is related to the breakdown voltage by:
\begin{resultbox}{Third-Order Intercept (IP3)}
% [VALUE SUBSTITUTION 2026-08-03 -- executed under the ratified sector re-keying
%  above; FLAGGED, not silent. Prior line preserved per Rule 12]: this equation
%  read
%    "V_{IP3} = \sqrt{\frac{4}{3}}\; V_{yield} \approx 1.155 \times 43.65
%     \approx 50.4 \text{ kV}"
%  The 50.4 kV is STRUCK. The IP3 form V_IP3 = sqrt(4/3) V_br is unchanged; what
%  changes is WHICH voltage is the breakdown reference, and the sector re-keying
%  rules that to be V_snap. Recompute on src/ave/core/constants.py:
%    V_IP3 = sqrt(4/3) * V_SNAP = 1.1547005 * 510,998.95 V = 590,051 V ~ 590 kV.
%  THE 11.7x MOVE IS DISCLOSED EXPLICITLY: the new value is exactly 1/sqrt(alpha)
%  times the old, because V_SNAP / V_YIELD = 1/sqrt(alpha) = 11.70624 identically.
%  Cross-check both directions:
%    sqrt(4/3) * V_YIELD = 50,404.8 V -- reproduces the struck 50.4 kV to the
%      digit, confirming the struck value's ARITHMETIC was right and only its
%      KEYING was wrong;
%    50,404.8 V * 11.70624 = 590,051 V.
%  This is a VALUE SUBSTITUTION downstream of an adjudicated physics ruling, not
%  a transcription correction, and is recorded as such in the commit message and
%  in _orchestration/docket-entries/2026-08-03-mr-handoff-mechanical.md.
\begin{equation}
    V_{IP3} = \sqrt{\frac{4}{3}}\; V_{snap} \approx 1.1547 \times 511.0 \approx 590 \text{ kV}
    \label{eq:ip3}
\end{equation}
\end{resultbox}

The IM3 sideband power relative to the fundamental is:
\begin{equation}
    P_{IM3} = P_{fund} - 3(V_{IP3,\text{dBm}} - P_{fund,\text{dBm}})
    \label{eq:im3_power}
\end{equation}

Table~\ref{tab:imd_levels} evaluates the predicted IM3 amplitude at representative drive levels.

\begin{table}[H]
\centering
\caption{Predicted 3rd-order IMD sideband levels vs.\ fractional drive voltage. The IM3 products are measurable above $-80$~dBc (thermal noise floor for a cryogenic receiver).}
\label{tab:imd_levels}
\small
\begin{tabularx}{\textwidth}{@{}@{@{}}r r r r@{}}
\toprule
$V/V_{yield}$ & Drive (kV) & IM3 Level (dBc) & Measurable? \\
\midrule
0.01 & 0.44  & $-160$ & No  \\
0.10 & 4.37  & $-100$ & No  \\
0.30 & 13.10 & $-70$  & Marginal \\
0.50 & 21.83 & $-54$  & Yes \\
0.70 & 30.56 & $-40$  & Yes \\
0.90 & 39.29 & $-20$  & Strong \\
\bottomrule
\end{tabularx}
\end{table}

\subsection{QED Comparison}

The Euler-Heisenberg effective Lagrangian predicts light-by-light scattering with cross-section:
\begin{equation}
    \sigma_{EH} = \frac{973\, \alpha^4}{10125\, \pi} \left(\frac{\omega}{m_e c^2}\right)^{\!6} r_e^2 \sim 10^{-65} \text{ cm}^2 \quad (\text{at optical frequencies})
\end{equation}
This is $\sim 10^{40}$ times smaller than the AVE prediction at the same frequency, because QED treats the vacuum non-linearity as a perturbative loop correction ($\alpha^4$), while AVE treats it as a macroscopic classical saturation with a definite voltage threshold. The distinction is experimentally decisive: at optical frequencies, QED predicts undetectable photon-photon scattering, while AVE predicts a specific, amplitude-dependent IM3 tone spectrum that becomes measurable above $\sim 30\%$ of $V_{yield}$ ($\sim 13$ kV).

\subsection{Experimental Falsification Criterion}

The IMD test constitutes a direct, tabletop falsification of the AVE framework:

\begin{enumerate}
    \item \textbf{Drive configuration:} Two co-propagating laser beams at frequencies $f_1, f_2$ (separated by $\Delta f \sim 1$~GHz for spectral resolution) are focused to achieve a combined electric field exceeding $30\%$ of $V_{yield}/\ell_{node}$.
    \item \textbf{Detection:} A high-sensitivity heterodyne receiver at $2f_1 - f_2$ and $2f_2 - f_1$.
    \item \textbf{Null result:} If no IM3 products are detected above $-80$~dBc at $V/V_{yield} > 0.5$, the Axiom~4 saturation kernel is falsified.
    \item \textbf{Positive result:} Detection of IM3 products scaling as $V^3$ (cubic power law) below IP3 constitutes direct evidence for a macroscopic, non-perturbative vacuum non-linearity consistent with the AVE saturation threshold.
\end{enumerate}

\noindent The predicted cubic power-law scaling of IM3 amplitude with drive level is the unique AVE signature. QED predicts a sextic ($V^6$) scaling from the $\alpha^4$ loop correction. Measuring the exponent to within $\pm 0.5$ would definitively distinguish the two frameworks.

\section{Computational Solver Selection}
\label{sec:solver_selection}

The AVE physics engine provides two distinct electromagnetic solvers, each axiomatically valid but optimised for different simulation domains. Selecting the appropriate solver is an engineering decision governed by the physical observable under study.

\subsection{Cartesian FDTD (Yee Grid)}

The 3D Finite-Difference Time-Domain engine (\texttt{ave.core.fdtd\_3d}) discretises Maxwell's equations on a standard Cartesian Yee staggered grid. Each cell is a lumped LC element (Axiom~1): the E-field update uses $\varepsilon_{eff}$ from the Axiom~4 saturation kernel, and the H-field update uses $\mu_{eff}$. The solver enforces the CFL stability condition $\Delta t \leq \Delta x / (c\sqrt{3})$.

\paragraph{Strengths.} The Yee grid is computationally efficient ($O(N)$ per timestep), trivially parallelisable, and well-suited to macroscopic engineering simulations where the simulation cell size $\Delta x \gg \ell_{node}$. At these scales, the intrinsic chirality of the vacuum lattice averages out over each computational cell, and the Cartesian discretisation introduces no systematic error in chirality-blind observables (energy density, ponderomotive force, impedance matching, S-parameters).

\paragraph{Limitations.} The Yee grid cannot resolve vacuum chirality. The K4 diamond topology of the physical lattice is approximated as an isotropic continuum. Observables that depend on the handedness of the vacuum (e.g., chiral acoustic emission, OAM coupling, topological Faraday rotation) are structurally invisible to the Cartesian stencil.

\subsection{K4-TLM (Tetrahedral Diamond Lattice)}

\noindent\textbf{Carrier note (2026-08-02, D1 ratification --- Grant 2026-07-03; solver descriptions below are unchanged).} The \texttt{k4\_tlm} engine described here implements the \textbf{degree-4 ($z = 4$) bipartite-FCC diamond} lattice (achiral $Fd\bar{3}m$). That object is the \emph{historical} engine substrate and is \textbf{re-tagged a non-canonical instrument}: the ratified \textbf{production carrier} named ``K4'' is the \textbf{chiral \emph{srs} $z = 3$ net} (the Sunada-K4 / Laves net that Axiom~1 names), and the migration is chartered. Read this subsection --- and the ``K4 diamond topology of the physical lattice'' phrase in the FDTD limitations above --- as describing \emph{an instrument}, not the vacuum's carrier. \emph{Registered open, not adjudicated here:} whether the port-handed \emph{achiral} diamond genuinely realises Axiom~1's \emph{chiral} $I4_132$ space group is an \textbf{open} question (the diamond's handedness is a dynamical $\kappa_{chiral}$ order parameter on an achiral lattice); the chirality claims in this subsection are left standing under that flag. Canonical: \kbleaf{ave-kb/common/vocabulary-register.md} (\texttt{def-4b1a2c} referents (a)/(b); \texttt{def-7c3f9e} chirality, status \emph{ambiguous}).

% [2026-08-02 CARRIER VOCABULARY correction (D1 ratification, Grant 2026-07-03). Prior wording
%  (git trail): "discretises the vacuum on the native diamond graph ($K_4$ complete graph)". The
%  word "native" asserted the degree-4 diamond as the vacuum's carrier. Authority:
%  ave-kb/common/vocabulary-register.md:526 (def-4b1a2c) -- the PRODUCTION CARRIER named "K4" is
%  the chiral srs z=3 net, RATIFIED by Grant 2026-07-03 (D1); and :534 referent (b) -- the engine
%  "K4" (src/ave/core/k4_tlm.py) is the degree-4 bipartite-FCC DIAMOND (achiral Fd-3m),
%  "re-tagged a non-canonical instrument by the 2026-07-03 D1 ratification". The solver
%  DESCRIPTION is factually accurate about the code and is retained unchanged; only the carrier
%  assertion is corrected. Same-file concordance: the branch-vertex paragraph already prints
%  "the chiral srs vacuum-net vertex (coordination z=3)".]
The K4 Transmission Line Matrix engine (\texttt{ave.core.k4\_tlm}) discretises the vacuum on a degree-4 bipartite diamond graph ($K_4$ complete graph). Each node has exactly four ports connected to its tetrahedral neighbours, with native chirality encoded in the scattering matrix. The solver applies the Axiom~4 saturation operator directly to the 4-port scatter coefficients.

\paragraph{Strengths.} The K4-TLM is the only solver that preserves the vacuum's topological chirality at the simulation level. It is required for any computation where the observable depends on the handedness of the lattice: chiral antenna design (torus knot coupling), gravitational frame-dragging (Lense--Thirring), and topological symmetry-breaking experiments.

\paragraph{Limitations.} The tetrahedral mesh is approximately $4\times$ more expensive per node than the Cartesian grid, and the non-orthogonal geometry complicates boundary conditions and source injection. For bulk electromagnetic observables (topological stress, energy, S-parameters), it produces identical results to FDTD at a higher computational cost.

\subsection{Selection Criteria}

Table~\ref{tab:solver_selection} provides the engineering decision matrix.

\begin{table}[H]
\centering
\caption{Solver selection guide for AVE simulations. Use FDTD for bulk engineering; use K4-TLM only when chirality is the observable.}
\label{tab:solver_selection}
\small
\begin{tabularx}{\textwidth}{@{}@{@{}}l c c@{}}
\toprule
\textbf{Observable / Application} & \textbf{FDTD} & \textbf{K4-TLM} \\
\midrule
Non-linear dielectric stress    & \checkmark & --- \\
Energy density \& impedance matching     & \checkmark & --- \\
Acoustic shear tensor ($\tau_{zx}$)     & \checkmark & --- \\
IMD spectroscopy                         & \checkmark & --- \\
Warp metric CFD (superluminal transit)   & \checkmark & --- \\
Chiral antenna coupling        & ---        & \checkmark \\
Gravitational lensing \& frame-dragging  & ---        & \checkmark \\
Topological Faraday rotation             & ---        & \checkmark \\
Vacuum birefringence (chirality-split)   & ---        & \checkmark \\
\bottomrule
\end{tabularx}
\end{table}

\paragraph{Boundary Conditions.} Both solvers support Mur 1st-order absorbing boundaries (default) and PML (Perfectly Matched Layer) polynomial-graded absorbing boundaries (\texttt{use\_pml=True}). PML is recommended for simulations requiring low spurious reflection ($<$\,60~dB) at the grid boundary, such as far-field antenna patterns and long-duration energy tracking. The PML implementation uses cubic-graded conductivity: $\sigma(d) = \sigma_{max}(d/d_{pml})^3$, with $\sigma_{max} = (m+1)/(150\pi\,\Delta x)$.

\paragraph{Default Yield Threshold.} Both FDTD and K4-TLM default to $V_{yield} = \sqrt{\alpha}\,V_{snap} \approx 43.65$~kV as the Axiom~4 nonlinear onset. Subatomic-scale simulations (e.g., bond energy solvers, Yang--Mills confinement) should override with \texttt{v\_yield=V\_SNAP} ($\approx 511$~kV).

\begin{summarybox}
\begin{itemize}
    \item This chapter demonstrates the direct engineering application of continuous Vacuum Circuit Analysis networks functioning natively within the AVE transmission line.
    \item No empirical fitting variables or Standard Model mass parameters are required; hardware geometry traces explicitly to the four macroscopic yielding limits.
    \item Hardware operation is verified by predicting exact reflection coefficients ($\Gamma$) and local lattice resonance boundaries in purely geometric terms.
\end{itemize}
\end{summarybox}

\begin{exercisebox}
\begin{enumerate}
    \item Derive the hardware fracture scale matching the rigorous mathematical limits of Vacuum Circuit Analysis when local strain exceeds the Axiom 4 topological dielectric threshold ($\varepsilon > 1.0$).
    \item Construct a native LC structural tensor illustrating how macroscopic power cascading forces the spatial substrate into localized saturation during active Vacuum Circuit Analysis operation.
\end{enumerate}
\end{exercisebox}
